\section{AIME I 2021}
        \begin{enumerate}
        	\item (\textit{AIME I 2021}) Zou and Chou are practicing their $100$-meter sprints by running $6$ races against each other. Zou wins the first race, and after that, the probability that one of them wins a race is $\frac23$ if they won the previous race but only $\frac13$ if they lost the previous race. The probability that Zou will win exactly $5$ of the $6$ races is $\frac mn$, where $m$ and $n$ are relatively prime positive integers. Find $m+n$.


 
 \textbf{Answer: }097



\textbf{Solution 1 (Casework 🪅)}

For the next five races, Zou wins four and loses one. Let $W$ and $L$ denote a win and a loss, respectively. There are five possible outcome sequences for Zou:

 <ol style="margin-left: 1.5em;"> \item $LWWWW$ 

 \item $WLWWW$ 

 \item $WWLWW$ 

 \item $WWWLW$ 

 \item $WWWWL$

 </ol>We proceed with casework:

 \textbf{<u>Case (1): Sequences \#1-4, in which Zou does not lose the last race.</u>}

 The probability that Zou loses a race is $\frac13,$ and the probability that Zou wins the next race is $\frac13.$ For each of the three other races, the probability that Zou wins is $\frac23.$

 There are four sequences in this case. The probability of one such sequence is $\left(\frac13\right)^2\left(\frac23\right)^3.$

 \textbf{<u>Case (2): Sequence \#5, in which Zou loses the last race.</u>}

 The probability that Zou loses a race is $\frac13.$ For each of the four other races, the probability that Zou wins is $\frac23.$

 There is one sequence in this case. The probability is $\left(\frac13\right)^1\left(\frac23\right)^4.$

 \textbf{<u>Answer</u>}

 The requested probability is 
 \[4\left(\frac13\right)^2\left(\frac23\right)^3+\left(\frac13\right)^1\left(\frac23\right)^4=\frac{32}{243}+\frac{16}{243}=\frac{48}{243}=\frac{16}{81},\] from which the answer is $16+81=\boxed{097}.$

 





\textbf{Solution 2 (Casework but Bashier)}

We have $5$ cases, depending on which race Zou lost. Let $\text{W}$ denote a won race, and $\text{L}$ denote a lost race for Zou. The possible cases are $\text{WWWWL, WWWLW, WWLWW, WLWWW, LWWWW}$. The first case has probability $\left(\frac{2}{3} \right)^4 \cdot \frac{1}{3} = \frac{16}{3^5}$. The second case has probability $\left( \frac{2}{3} \right)^3 \cdot \frac{1}{3} \cdot \frac{1}{3} = \frac{8}{3^5}$. The third has probability $\left( \frac{2}{3} \right)^2 \cdot \frac{1}{3} \cdot \frac{1}{3} \cdot \frac{2}{3} = \frac{8}{3^5}$. The fourth has probability $\frac{2}{3} \cdot \frac{1}{3} \cdot \frac{1}{3} \cdot \left( \frac{2}{3} \right)^2 = \frac{8}{3^5}$. Lastly, the fifth has probability $\frac{1}{3} \cdot \frac{1}{3} \cdot \left( \frac{2}{3} \right)^3 = \frac{8}{3^5}$. Adding these up, the total probability is $\frac{16 + 8 \cdot 4}{3^5} = \frac{16 \cdot 3}{3^5} = \frac{16}{81}$, so $m+n = \boxed{097}$.

 This is for if you're paranoid like me, and like to sometimes write out all of the cases if there are only a few.

 





\textbf{Solution 3 (Even More Casework)}

\textbf{ Case 1: Zou loses the second race }

 In this case, Zou must win the rest of the races. Thus, our probability is $\frac{8}{243}$.

 \textbf{ Case 2: Zou loses the last race }

 There is only one possibility for this, so our probability is $\frac{16}{243}$.

 \textbf{ Case 3: Neither happens }

 There are three ways that this happens. Each has one loss that is not the last race. Therefore, the probability that one happens is $\frac{1}{3} \cdot \frac{1}{3} \cdot \frac{2}{3} \cdot \frac{2}{3} \cdot \frac{2}{3} = \frac{8}{243}$. Thus, the total probability is $\frac{8}{243} \cdot 3 = \frac{24}{243}$.

 Adding these up, we get $\frac{48}{243} = \frac{16}{81}$, so $16+81=\boxed{097}$.

 





\textbf{Solution 4 (Observations)}

Note that Zou wins one race. The probability that he wins the last race is $\left(\frac{2}{3}\right)^4\left(\frac{1}{3}\right)=\frac{16}{243}.$ Now, if he doesn't win the last race, then there must be two races where the winner of the previous race loses. We can choose any $4$ of the middle races for Zou to win. So the probability for this case is $4\left(\frac{2}{3}\right)^3\left(\frac{1}{3}\right)^2=\frac{32}{243}.$ Thus, the answer is $\frac{16}{243}+\frac{32}{243}=\frac{16}{81}\implies\boxed{097}.$

 




	\item (\textit{AIME I 2021}) In the diagram below, $ABCD$ is a rectangle with side lengths $AB=3$ and $BC=11$, and $AECF$ is a rectangle with side lengths $AF=7$ and $FC=9,$ as shown. The area of the shaded region common to the interiors of both rectangles is $\frac mn$, where $m$ and $n$ are relatively prime positive integers. Find $m+n$.


 
                    \begin{figure}[H]
                        \centering
                        \includegraphics[width=0.5\textwidth]{Problem/AIME_2021_I_2_Problem_Figure_1.png}
                    \end{figure}
                    


 
 \textbf{Answer: }109



\textbf{Solution 1 (Similar Triangles)}

Let $G$ be the intersection of $AD$ and $FC$.From vertical angles, we know that $\angle FGA= \angle DGC$. Also, because we are given that $ABCD$ and $AFCE$ are rectangles, we know that $\angle AFG= \angle CDG=90 ^{\circ}$.Therefore, by AA similarity, we know that $\triangle AFG\sim\triangle CDG$.

 Let $AG=x$. Then, we have $DG=11-x$. By similar triangles, we know that $FG=\frac{7}{3}(11-x)$ and $CG=\frac{3}{7}x$. We have $\frac{7}{3}(11-x)+\frac{3}{7}x=FC=9$.

 Solving for $x$, we have $x=\frac{35}{4}$.The area of the shaded region is just $3\cdot \frac{35}{4}=\frac{105}{4}$.

 Thus, the answer is $105+4=\framebox{109}$. 

 





\textbf{Solution 2 (Similar Triangles)}

Again, let the intersection of $AE$ and $BC$ be $G$. By AA similarity, $\triangle AFG \sim \triangle CDG$ with a $\frac{7}{3}$ ratio. Define $x$ as $\frac{[CDG]}{9}$. Because of similar triangles, $[AFG] = 49x$. Using $ABCD$, the area of the parallelogram is $33-18x$. Using $AECF$, the area of the parallelogram is $63-98x$. These equations are equal, so we can solve for $x$ and obtain $x = \frac{3}{8}$. Thus, $18x = \frac{27}{4}$, so the area of the parallelogram is $33 - \frac{27}{4} = \frac{105}{4}$.

 Finally, the answer is $105+4=\boxed{109}$.

 





\textbf{Solution 3 (Pythagorean Theorem)}

Let the intersection of $AE$ and $BC$ be $G$, and let $BG=x$, so $CG=11-x$. 

 By the Pythagorean theorem, ${AG}^2={AB}^2+{BG}^2$, so $AG=\sqrt{x^2+9}$, and thus $EG=9-\sqrt{x^2+9}$. 

 By the Pythagorean theorem again, ${CG}^2={EG}^2+{CE}^2$: 
 \[11-x=\sqrt{7^2+(9-\sqrt{x^2+9})^2}.\]

 Solving, we get $x=\frac{9}{4}$, so the area of the parallelogram is $3\cdot\left(11-\frac{9}{4}\right)=\frac{105}{4}$, and $105+4=\framebox{109}$.

 





\textbf{Solution 4 (Pythagorean Theorem)}

Let $P = AD \cap FC$, and $K = AE \cap BC$. Also let $AP = x$.

 $CK$ also has to be $x$ by parallelogram properties. Then $PD$ and $BK$ must be $11-x$ because the sum of the segments has to be $11$.

 We can easily solve for $PC$ by the Pythagorean Theorem:
 \begin{align*} DC^2 + PD^2 &= PC^2\\ 9 + (11-x)^2 &= PC^2 \end{align*}
It follows shortly that $PC = \sqrt{x^2-22x+30}$.

 Also, $FC = 9$, and $FP + PC = 9$. We can then say that $PC = \sqrt{x^2-22x+30}$, so $FP = 9 - \sqrt{x^2-22x+30}$.

 Now we can apply the Pythagorean Theorem to $\triangle AFP$.
 \begin{align*} AF^2 + FP^2 = AP^2\\ 49 + \left(9 - \sqrt{x^2-22x+30}\right)^2 = x^2  \end{align*}


 This simplifies (not-as-shortly) to $x = \dfrac{35}{4}$. Now we have to solve for the area of $APCK$. We know that the height is $3$ because the height of the parallelogram is the same as the height of the smaller rectangle.

 From the area of a parallelogram (we know that the base is $\dfrac{35}{4}$ and the height is $3$), it is clear that the area is $\dfrac{105}{4}$, giving an answer of $\boxed{109}$. 

 





\textbf{Solution 5 (Coordinate Geometry)}

Suppose $B=(0,0).$ It follows that $A=(0,3),C=(11,0),$ and $D=(11,3).$

 Since $AECF$ is a rectangle, we have $AE=FC=9$ and $EC=AF=7.$ The equation of the circle with center $A$ and radius $\overline{AE}$ is $x^2+(y-3)^2=81,$ and the equation of the circle with center $C$ and radius $\overline{CE}$ is $(x-11)^2+y^2=49.$

 We now have a system of two equations with two variables. Expanding and rearranging respectively give
 \begin{align*} x^2+y^2-6y&=72, &(1) \\ x^2+y^2-22x&=-72. &(2) \end{align*}
Subtracting $(2)$ from $(1),$ we obtain $22x-6y=144.$ Simplifying and rearranging produce 
 \[x=\frac{3y+72}{11}. \hspace{34.5mm} (*)\]Substituting $(*)$ into $(1)$ gives 
 \[\left(\frac{3y+72}{11}\right)^2+y^2-6y=72,\] which is a quadratic of $y.$ We clear fractions by multiplying both sides by $11^2=121,$ then solve by factoring:
 \begin{align*} \left(3y+72\right)^2+121y^2-726y&=8712 \\ \left(9y^2+432y+5184\right)+121y^2-726y&=8712 \\ 130y^2-294y-3528&=0 \\ 2(5y+21)(13y-84)&=0 \\ y&=-\frac{21}{5},\frac{84}{13}. \end{align*}
Since $E$ is in Quadrant IV, we have $E=\left(\frac{3\left(-\frac{21}{5}\right)+72}{11},-\frac{21}{5}\right)=\left(\frac{27}{5},-\frac{21}{5}\right).$ It follows that the equation of $\overleftrightarrow{AE}$ is $y=-\frac{4}{3}x+3.$ 

 Let $G$ be the intersection of $\overline{AD}$ and $\overline{FC},$ and $H$ be the intersection of $\overline{AE}$ and $\overline{BC}.$ Since $H$ is the $x$-intercept of $\overleftrightarrow{AE},$ we get $H=\left(\frac94,0\right).$

 By symmetry, quadrilateral $AGCH$ is a parallelogram. Its area is $HC\cdot AB=\left(11-\frac94\right)\cdot3=\frac{105}{4},$ from which the requested sum is $105+4=\boxed{109}.$ 

 





\textbf{Solution 6 (Trigonometry)}

Let the intersection of $AE$ and $BC$ be $G$. It is useful to find $\tan(\angle DAE)$, because  $\tan(\angle DAE)=\frac{3}{BG}$  and  $\frac{3}{\tan(\angle DAE)}=BG$. From there, subtracting the areas of the two triangles from the larger rectangle, we get  Area = $33-3BG=33-\frac{9}{\tan(\angle DAE)}$.

 let $\angle CAD = \alpha$. Let $\angle CAE = \beta$. Note, $\alpha+\beta=\angle DAE$. 

 $\alpha=\tan^{-1}\left(\frac{3}{11}\right)$

 $\beta=\tan^{-1}\left(\frac{7}{9}\right)$

 $\tan(\angle DAE) = \tan\left(\tan^{-1}\left(\frac{3}{11}\right)+\tan^{-1}\left(\frac{7}{9}\right)\right) = \frac{\frac{3}{11}+\frac{7}{9}}{1-\frac{3}{11}\cdot\frac{7}{9}} = \frac{\frac{104}{99}}{\frac{78}{99}} = \frac{4}{3}$  

 $\mathrm{Area}=33-\frac{9}{\frac{4}{3}} = 33-\frac{27}{4 } = \frac{105}{4}$. The answer is $105+4=\boxed{109}$.

 





\textbf{Solution 7}

A self-admittedly worse solution 1

 Construct a point at line $\overline{BC}$, and name it $G$.Let $a$ = $\overline{AG}$Let $b$ = $\overline{GE}$Let $c$ = $\overline{BG}$Let $d$ = $\overline{AB}$Let $e$ = $\overline{GC}$Let $f$ = $\overline{AF}$

 We know $\overline{AB} = 3$, so let $d$ = $3$. We also know that $\overline{AF} = 7$, so let $f$ = $7$.

 Triangle $\triangle ABG \sim \triangle ECG$ due to all the angles being equal from alternate interior angles at point G and both of the triangles having right angles. Dividing $d$ and $f$ yields the ratio of the side lengths to be $\frac{3}{7}$

 We will now create several equations with these variables.$d = 3$$c + e = 11$$f =  7$$a + b = 9$$\frac{d}{f} = \frac{3}{7}$$\frac{a}{e} = \frac{3}{7}$$\frac{c}{b} = \frac{3}{7}$

 We can find $b$ to be $9 - a$ from the fourth equation, and then find $c$ to be $7c = 3b$ from the seventh equation. Then substitute $b$ for $9 - a$ to find that $c = \frac{3}{7}\left(9 - a \right)$. Now we just substitute $c$ into our second equation, $c + e = 11$ to yield that $e$ is $e = 11 - \frac{3}{7}\left(9 - a \right)$. That's not useful right now, but will be once we look at our sixth equation, which shows that $\frac{a}{e} = \frac{3}{7}$, and we yield that $e$ is equal to $3e = 7a$ which shows that when simplified we get $e = \frac{7a}{3}$. Now we can see that we have 2 equations with 2 variables. When substituting $e$ into the first equation, which is $\frac{a}{e} = \frac{3}{7}$, we get 
 \[\frac{7a}{3} = 11 - \frac{3}{7}\left(9 - a \right)\] We then yield that $15 = 4a$, which gives us $a = \frac{15}{4}$. We then substitute $a$ into our fourth equation, which is $a + b = 9$, and then we get that $b = 9 - \frac{15}{4}$, and then we get that $b$ is equal to $\frac{21}{4}$. Then we use our seventh equation, which is $\frac{c}{b} = \frac{3}{7}$. We then get $7c = 3b$. We substitute $b$ in to get that $c$ is $\frac{9}{4}$ Then we use our second equation which is $c + e = 11$ and then we substitute $c$ into it to get that $e$ is $11 - \frac{9}{4}$. Then we get that $e = \frac{35}{4}$. Then to find the area we multiply $e$ by the height, $d$ to get $\frac{35}{4} \cdot 3$ to get $\frac{105}{4}$, which is $\frac{m}{n}$, and we get that $m = 105$ and $n = 4$, so 
 \[m + n = 105 + 4 = \boxed{109}\]

 





\textbf{Solution 8(Extremely simple solution)}

Let $EC \cap AD = P$ and $AE \cap BC = Q$. By $AA$ similarity, $\triangle PDC \sim \triangle AFP$ with corresponding sides $CD$ and $AF$. Let $[AFP] = [QEC] = X$ and $[ABQ] = [CDP] = Y$. Let the shaded parallelogram have area $M$. We know$[AFCE] = [AFP] + M + [QEC] = 2X + M = 7 \cdot 9 = 63$ and $[ABCD] = [ABQ] + M + CDP = 2Y + M = 3 \cdot 11 = 33$ thus subtracting them gives $2X - 2Y = 63 - 33 \implies X - Y = 15 \implies Y = X - 15$. Now note $\frac{[CDP]}{[AFP]} = (\frac{CD}{AF})^2 = (\frac{3}{7})^2 = \frac{9}{49}$ but $[CDP] = Y = X - 15$ and $[AFP] = X$ thus we solve $\frac{X - 15}{X} = \frac{9}{49} \implies 49X - 49 \cdot 15 = 9X \implies 40X = 49 \cdot 15 \implies 8X = 49 \cdot 3 \implies X = \frac{49 \cdot 3}{8}$. We want $M$. However, note $2X + M = 63 \implies M = 63 - 2X \implies M = 63 - 2(\frac{49 \cdot 3}{8}) = 63 - \frac{49 \cdot 3}{4} = \frac{63 \cdot 4 - 49 \cdot 3}{4} = \frac{7(9 \cdot 4 - 7 \cdot 3)}{4} = \frac{7(36 - 21)}{4} = \frac{7 \cdot 15}{4} = \frac{105}{4} \implies \boxed{109}$

 




	\item (\textit{AIME I 2021}) Find the number of positive integers less than $1000$ that can be expressed as the difference of two integral powers of $2.$


 
 \textbf{Answer: }050



\textbf{Solution 1}

We want to find the number of positive integers $n<1000$ which can be written in the form $n = 2^a - 2^b$ for some non-negative integers $a > b \ge 0$ (note that if $a=b$, then $2^a-2^b = 0$). We first observe $a$ must be at most 10; if $a \ge 11$, then $2^a - 2^b \ge 2^{10} > 1000$. As $2^{10} = 1024 \approx 1000$, we can first choose two different numbers $a > b$ from the set $\{0,1,2,\ldots,10\}$ in $\binom{11}{2}=55$ ways. This includes $(a,b) = (10,0)$, $(10,1)$, $(10,2)$, $(10,3)$, $(10,4)$ which are invalid as $2^a - 2^b > 1000$ in this case. For all other choices $a$ and $b$, the value of $2^a - 2^b$ is less than 1000.

 We claim that for all other choices of $a$ and $b$, the values of $2^a - 2^b$ are pairwise distinct. More specifically, if $(a_1,b_1) \neq (a_2,b_2)$ where $10 \ge a_1 > b_1 \ge 0$ and $10 \ge a_2 > b_2 \ge 0$, we must show that $2^{a_1}-2^{b_1} \neq 2^{a_2} - 2^{b_2}$. Suppose otherwise for sake of contradiction; rearranging yields $2^{a_1}+2^{b_2} = 2^{a_2}+2^{b_1}$. We use the fact that every positive integer has a unique binary representation:

 If $a_1 \neq b_2$ then $\{a_1,b_2\} = \{a_2,b_1\}$; from here we can deduce either $a_1=a_2$ and $b_1=b_2$ (contradicting the assumption that $(a_1,b_1) \neq (a_2,b_2)$, or $a_1=b_1$ and $a_2=b_2$ (contradicting the assumption $a_1>b_1$ and $a_2>b_2$).

 If $a_1 = b_2$ then $2^{a_1}+2^{b_2} = 2 \times 2^{a_1}$, and it follows that $a_1=a_2=b_1=b_2$, also contradicting the assumption $(a_1,b_1) \neq (a_2,b_2)$. Hence we obtain contradiction.*

 Then there are $\binom{11}{2}-5$ choices for $(a,b)$ for which $2^a - 2^b$ is a positive integer less than 1000; by the above claim, each choice of $(a,b)$ results in a different positive integer $n$. Then there are $55-5 = \boxed{050}$ integers which can be expressed as a difference of two powers of 2.

 


 
                \begin{itemize}
                    \item Note: The uniqueness of binary representation could be rather easily proven, but if you cannot convince yourself on the spot that this is the case, consider the following alternative proof. Let $(a_1,b_1) \neq (a_2,b_2)$ where $10 \ge a_1 > b_1 \ge 0$ and $10 \ge a_2 > b_2 \ge 0$ and $2^{a_1}-2^{b_1} = 2^{a_2} - 2^{b_2}$, for the sake of contradiction. Therefore $\deg_{2}(2^{a_1}-2^{b_1})=\deg_{2}(2^{a_2}-2^{b_2})$, or $b_1=b_2$. Plugging in, we see that $2^{a_1}=2^{a_2}$, or $a_1=a_2$, contradiction.
                \end{itemize}


                Note by Ross Gao

 





\textbf{Solution 2 (Casework)}

\textbf{Case 1:} When our answer is in the form $2^n-2^i$, where $i$ is an integer such that $0\le i\le 4$. 

 We start with the subcase where it is $2^n-2^0$, for some integer $n$ where $n>0$ (this is because the case where $n=0$ yields $2^0-2^0=0$, which doesn't work because it must be a positive integer.) Note that $2^{10}=1024$, and $2^9=512$. Our answer needs to be less than $1000$, so the maximum possible result (in this case) is $2^9-2^0$. Our lowest result is $2^1-2^0$. All the positive powers of two less than $1024$ work, so we have $9$ possibilities for this subcase. For subcases $i=1, i=2, i=3,$ and $i=4$, we have $8, 7, 6,$ and $5$ possibilities, respectively.

 \textbf{Case 2:} When our answer is in the form of $2^n-2^j$, where $j$ is an integer such that $5\le j\le 9$.

 We can start with the subcase where $j=5$. We notice that $2^5=32$, and $2^{10}-2^5=992$ which is less than $1000$, so the greatest result in this subcase is actually $2^{10}-2^5$, and the lowest is $2^6-2^5$. Thus, we have $5$ possibilities. For the other four subcases, we have $4, 3, 2,$ and $1$ possibilities, respectively.

 \textbf{Answer:}We note that these are our only cases, as numbers in the form of $2^n-2^{10}$ and beyond are greater than $1000$. 

 Thus, our result is $9+8+7+6+5+5+4+3+2+1=(9+8+7+6+5+4+3+2+1)+5=\boxed{50}$. ~jehu26

 





\textbf{Solution 3 (Bash)}

We look for all positive integers of the form $2^a-2^b<1000,$ where $0\leq b<a.$ Performing casework on $a,$ we can enumerate all possibilities in the table below:
 \[\begin{array}{c|c} & \\ [-2.25ex] \boldsymbol{a} & \boldsymbol{b} \\  \hline  & \\ [-2ex] 1 & 0 \\   2 & 0,1 \\ 3 & 0,1,2 \\ 4 & 0,1,2,3 \\ 5 & 0,1,2,3,4 \\ 6 & 0,1,2,3,4,5 \\ 7 & 0,1,2,3,4,5,6 \\ 8 & 0,1,2,3,4,5,6,7 \\ 9 & 0,1,2,3,4,5,6,7,8 \\ 10 & \xcancel{0},\xcancel{1},\xcancel{2},\xcancel{3},\xcancel{4},5,6,7,8,9 \\ [0.5ex] \end{array}\]As indicated by the X-marks, the ordered pairs $(a,b)=(10,0),(10,1),(10,2),(10,3),(10,4)$ generate $2^a-2^b>1000,$ which are invalid.

 \textbf{\textit{Note that each of the remaining ordered pairs generates one unique desired positive integer.}} 

 We prove this statement as follows:

 <ol style="margin-left: 1.5em;"> \item The positive integers generated for each value of $a$ are clearly different. 

 \item For all integers $k$ such that $1\leq k\leq9,$ the largest positive integer generated for $a=k$ is $1$ less than the smallest positive integer generated for $a=k+1.$

 </ol>Together, we have justified our claim in bold. The answer is 
 \[1+2+3+4+5+6+7+8+9+5=\boxed{050}.\]

 





\textbf{Solution 4 (Faster Way)}

Because the difference is less than $1000$, we can simply list out all numbers that satisfy $2^n < 1000$. We get $0 \le n < 10$, where n is an integer. Because the sequence $2^n$ is geometric, the difference of any two terms will be unique. $\binom{10}{2}$ will be the number of differences for $0\le n < 10$. However, we also need to consider the case in which $n=10$. With simple counting, we find that $5$ numbers: $(32, 64, 128, 256, 512)$ could be subtracted from $1024$, which makes another 5 cases. There is no need to check for higher exponents since the lowest difference would be $2^{11} - 2^{10} = 1024$, which exceeds $1000$.

 Thus, the final answer is $\binom{10}{2} + 5 = \boxed{050}.$

 





\textbf{Solution 5 (Faster than most other solutions, Come and Go)}

Note $2^{11} - 2^{10} = 2^{10} \cdot 1 = 2^{10} > 1000 \therefore 2^a - 2^b :~ a > b~ \text{and}~ 0 \le a,b \le 10$. For the extrema $a = 10$, $2^b > 24 \implies 5 \ge b \ge 9$. For all other cases $a \neq 10$, $\#(B) = a$, where $\#(S)$ denotes the number of plausible $S$, or cardinality of $S$. Then, for $a \neq 10 \implies ~ \text{total number of possible}~b = 45$ and $a = 10 \implies ~ \text{total number of possible}~b = 5$. The count is $45+5 = \boxed{050}$.

 




	\item (\textit{AIME I 2021}) Find the number of ways $66$ identical coins can be separated into three nonempty piles so that there are fewer coins in the first pile than in the second pile and fewer coins in the second pile than in the third pile.


 
 \textbf{Answer: }331



\textbf{Solution 1}

Suppose we have $1$ coin in the first pile. Then $(1, 2, 63), (1, 3, 62), \ldots, (1, 32, 33)$ all work for a total of $31$ piles. Suppose we have $2$ coins in the first pile, then $(2, 3, 61), (2, 4, 60), \ldots, (2, 31, 33)$ all work, for a total of $29$. Continuing this pattern until $21$ coins in the first pile, we have the sum 
 \begin{align*} 31+29+28+26+25+\cdots+4+2+1 &= (31+28+25+22+\cdots+1)+(29+26+23+\cdots+2) \\ &= 176+155 \\ &= \boxed{331}. \end{align*}


 





\textbf{Solution 2}

We make an equation: $a+b+c=66,$ where $a<b<c.$ We don't have a clear solution, so we'll try complementary counting. First, let's find where $a\geq b\geq c.$ By stars and bars, we have $\dbinom{65}{2}=2080$ to assign positive integer solutions to $a + b + c = 66.$ Now we need to subtract off the cases where it doesn't satisfy the condition.

 We start with $a = b.$ We can write that as $2b + c = 66.$ We can find there are 32 integer solutions to this equation. There are $32$ solutions for $b=c$ and $a = c$ by symmetry. We also need to add back $2$ because we subtracted $(a,b,c)=(22,22,22)$ $3$ times.

 We then have to divide by $6$ because there are $3!=6$ ways to order $a, b,$ and $c.$ Therefore, we have $\dfrac{\dbinom{65}{2}-96+2}{6} = \dfrac{1986}{6} = \boxed{331}.$

 





\textbf{Solution 3}

Let the piles have $a, b$ and $c$ coins, with $0 < a < b < c$. Then, let $b = a + k_1$, and $c = b + k_2$, such that each $k_i \geq 1$. The sum is then $a + a+k_1 + a+k_1+k_2 = 66 \implies 3a+2k_1 + k_2 = 66$. This is simply the number of positive solutions to the equation $3x+2y+z = 66$. Now, we take cases on $a$. 

 If $a = 1$, then $2k_1 + k_2 = 63 \implies 1 \leq k_1 \leq 31$. Each value of $k_1$ corresponds to a unique value of $k_2$, so there are $31$ solutions in this case. Similarly, if $a = 2$, then $2k_1 + k_2 = 60 \implies 1 \leq k_1 \leq 29$, for a total of $29$ solutions in this case. If $a = 3$, then $2k_1 + k_2 = 57 \implies 1 \leq k_1 \leq 28$, for a total of $28$ solutions. In general, the number of solutions is just all the numbers that aren't a multiple of $3$, that are less than or equal to $31$. 

 We then add our cases to get 
 \begin{align*} 1 + 2 + 4 + 5 + \cdots + 31 &= 1 + 2 + 3 + \cdots + 31 - 3(1 + 2 + 3 + \cdots + 10) \\ &= \frac{31(32)}{2} - 3(55) \\ &= 31 \cdot 16 - 165 \\ &= 496 - 165 \\ &= \boxed{331} \end{align*}
as our answer.

 





\textbf{Solution 4}

Let the first pile have $a$ coins, the second $a+b$ coins, and the third $a+b+c$ coins, where $a$, $b$, and $c$ are strictly positive integers. Thus the total number of coins is $3a+2b+c=66$. Perform the substitution $x=a-1$, $y=b-1$, and $z=c-1$ to yield the equation $3x+2y+z=60$, where $x$, $y$, and $z$ are instead nonnegative integers.

 From here we can set up the \href{https://artofproblemsolving.com/wiki/index.php?title=Generating\_function}{generating function} $(x^0+x^3+\cdots+x^{60})(x^0+x^2+\cdots+x^{60})(x^0+x^1+\cdots+x^{60})$. We need to find the coefficient of $x^{60}$. Multiplying the second and third polynomials with clever reasoning returns

 $(x^0+x^3+\cdots+x^{60})(x^0+x^1+2x^2+2x^3+\cdots+31x^{60}+\cdots+2x^{117}+2x^{118}+x^{119}+x^{120})$

 where in the second polynomial, for every two terms, the coefficient increases or decreases by one (depending on which side of the polynomial the term resides).

 One can notice here that for every term in the first polynomial there exists one and only one term in the second polynomial that, when multiplied, yield $x^{60}$. Furthermore, we need only consider the <u>coefficients</u> of the second polynomial.

 The corresponding coefficient for $x^{60}$ is $1$, for $x^{57}$ is $2$, and for $x^{54}$ is $4$. We notice the pattern: increase by one, increase by two, and so on. When does this pattern stop? For $x^0$, the corresponding coefficient is $31$, and we notice that $31\equiv1\mod3$. As a result, we know that the pattern has $21$ terms, and we can take advantage of the first $20$ by symmetry.

 The answer is simply $10(1+29)+31=\boxed{331}$.

 Solution 4

 





\textbf{Solution 5}

Suppose we start with an equal pile of coins with 21, 21, and 21. There is 3 extra coins we must distribute to the third and second piles.

 Thus, the number of piles that have 21 as the least amount of coins will come from the equation. $3 = x + y$ where x and y are the second and third piles respectively. However, the restriction to this equation is that x and y can not equal 0.

 Therefore, we can set $x' = x + 1$ and $y' = y + 1$. The equation then turns into $1 = x' + y'$ Using stars and bars, the number of solutions to this equation becomes $\binom{2}{1}$. However, we must divide by 2 since we want the solution where $x' < y'$

 Continuing this thinking for the other cases where the least pile can be 20, 19, etc, we get the sum:

 $\frac{\binom{2}{1}}{2}$ + $\frac{\binom{5}{1} - 1}{2}$ + $\frac{\binom{8}{1}}{2}$ ... + $\frac{\binom{62}{1}}{2}$

 Note: We subtract one for some of them because it represents the case where $x' = y'$ which is not allowed.

 This sum simplifies into an arithmetic sequence with a few 1's subtracted:

 $\frac{((2\cdot21 + 3\cdot(\frac{20*21}{2})) + -1(10))}{2} = \boxed{331}$

 
                \begin{itemize}
                    <li>
                \begin{itemize}
                    \item 
                \end{itemize}


                </li>
                \end{itemize}


                




	\item (\textit{AIME I 2021}) Call a three-term strictly increasing arithmetic sequence of integers special if the sum of the squares of the three terms equals the product of the middle term and the square of the common difference. Find the sum of the third terms of all special sequences.


 
 \textbf{Answer: }031



\textbf{Solution 1}

Let the terms be $a-b$, $a$, and $a+b$. Then we want $(a-b)^2+a^2+(a+b)^2=ab^2$, or $3a^2+2b^2=ab^2$. Rearranging, we get $b^2=\frac{3a^2}{a-2}$. Simplifying further, $b^2=3a+6+\frac{12}{a-2}$. Looking at this second equation, since the right side must be an integer, $a-2$ must equal $\pm1, 2, 3, 4, 6, 12$. Looking at the first equation, we see $a>2$ since $b^2$ is positive. This means we must test $a=3, 4, 5, 6, 8, 14$. After testing these, we see that only $a=5$ and $a=14$ work which give $b=5$ and $b=7$ respectively. Thus the answer is $10+21=\boxed{031}$.~JHawk0224

 Note: You could also use synthetic division to get $b^2=\frac{3a^2}{a-2}$.

 





\textbf{Solution 2}

Let the common difference be $d$ and let the middle term be $x$. Then, we have that the sequence is
 \[x-d,~x,~x+d.\]This means that the sum of the squares of the 3 terms of the sequence is 
 \[(x-d)^2+x^2+(x+d)^2=x^2-2xd+d^2+x^2+x^2+2xd+d^2=3x^2+2d^2.\]We know that this must be equal to $xd^2,$ so we can write that
 \[3x^2+2d^2=xd^2,\]and it follows that
 \[3x^2-xd^2+2d^2=3x^2-\left(d^2\right)x+2d^2=0.\]

 Now, we can treat $d$ as a constant and use the quadratic formula to get
 \[x=\frac{d^2\pm \sqrt{d^4-4(3)(2d^2)}}{6}.\]We can factor pull $d^2$ out of the square root to get
 \[x=\frac{d^2\pm d\sqrt{d^2-24}}{6}.\]Here, it is easy to figure out the values of $d$. Let $\sqrt{d^2-24} = k$, then $d^2-k^2=24$ which is $(d+k)(d-k)=24,$ note that $d$, $k$ are integers. Examining the parity, we find that $d+k$ and $d-k$ are of the same parity. Now, we solve by factoring. We can find that $d=5$ and $d=7$ are the only positive integer values of $d$ that make $\sqrt{d^2-24}$ a positive integer.$^{*}$ $d=5$ gives $x=5$ and $x=\frac{10}{3}$, but we can ignore the latter. $d=7$ gives $x=14$, as well as a fraction which we can ignore. 

 Since $d=5,~x=5$ and $d=7, x=14$ are the only two solutions and we want the sum of the third terms, our answer is $(5+5)+(7+14)=10+21=\boxed{031}$. -BorealBear, minor edit by Kinglogic

 
$^*$To prove this, let $\sqrt{d^2-24} = k$, then $d^2-k^2=24$ which is $(d+k)(d-k)=24,$ then remembering that $d$ and $k$ are integers see if you can figure it out. -\href{https://artofproblemsolving.com/wiki/index.php?title=User:PureSwag}{PureSwag}

 





\textbf{Solution 3}

Proceed as in solution 2, until we reach 
 \[3x^2+2d^2=xd^2,\] Write 

 $d^2=\frac{3x^2}{x-2}$, it follows that $x-2=3k^2$ for some (positive) integer k and $k \mid x$.

 Taking both sides modulo $k$, $-2 \equiv 0 \pmod{k}$, so $k \mid 2 \rightarrow k=1,2$. 

 When $k=1$, we have $x=5$ and $d=5$. When $k=2$, we have $x=14$ and $d=7$.Summing the two cases, we have $10+21=\boxed{031}$.

 -Ross Gao

 





\textbf{Solution 4 (Combining Solution 1 and Solution 3)}

As in Solution 1, write the three integers in the sequence as $a-d$, $a$, and $a+d$.

 Then the sum of the squares of the three integers is $(a-d)^2+a^2+(a+d)^2 = 3a^2+2d^2$.

 Setting this equal to the middle term times the common difference squared, which is $ad^2$, 

 and solving for $d^2$ we get:

 $3a^2+2d^2 = ad^2 \implies ad^2-2d^2 = 3a^2 \implies d^2(a-2) = 3a^2 \implies d^2 = \frac{3a^2}{a-2}$

 The numerator has to be positive, so the denominator has to be positive too for the sequence 

 to be strictly increasing; that is, $a>2$. 

 For $\frac{3a^2}{a-2}$ to be a perfect square, $\frac{3}{a-2}$ must be a perfect square as well.

 This means that $a-2$ is divisible by 3, and whatever left over is a perfect square.

 We can express this as an equation: let the perfect square left over be $n^2$. Then:

 $3n^2 = a-2$. Now when you divide the numerator and denominator by 3, you are left with

 $d^2 = \frac{a^2}{n^2} \implies d = \frac{a}{n}$. Because the sequence is of integers, d must also be an 

 integer, which means that $n$ must divide $a$. 

 Taking the above equation we can solve for $a$: $3n^2 = a-2 \implies a = 3n^2+2$. 

 This means that $3n^2+2$ is divisible by $n$. $3n^2$ is automatically divisible by $n$, so 

 $2$ must be divisible by $n$. Then $n$ must be either of $\{1,2\}$. Plugging back into the equation, 

 $n = 1 \implies a = 5 \implies d = 5$, so $a+d = 5+5 = 10$.

 $n = 2 \implies a = 14 \implies d = 7$, so $a+d = 14+7 = 21$.

 Finally, $10+21 = \boxed{031}$

 -KingRavi

 





\textbf{Solution 5}

Following from previous solutions, we derive $3x^2+2a^2=xa^2.$ We divide both sides to get $3\left(\frac{x}{a}\right)^2+2=x.$ Since $x$ is an integer, $\frac{x}{a}$ must also be an integer, so we have $x=pa$, for some factor $p$. We then get $3p^2+2=pa.$ We then take this to modulo $p$, getting $2\equiv 0 \pmod p.$ The only possibilities for $p$ are therefore 1 and 2. We plug these into $3p^2+2=pa$, for $a=5$ and $x=5$, giving us the sequence $0,5,10$, or $2a=14$ and $x=14$, for the sequence $7,14,21.$ $10+21 = \boxed{031}.$

 -RYang2

 




	\item (\textit{AIME I 2021}) Segments $\overline{AB}, \overline{AC},$ and $\overline{AD}$ are edges of a cube and $\overline{AG}$ is a diagonal through the center of the cube. Point $P$ satisfies $BP=60\sqrt{10}$, $CP=60\sqrt{5}$, $DP=120\sqrt{2}$, and $GP=36\sqrt{7}$. Find $AP.$


 
 \textbf{Answer: }192



\textbf{Solution 1}

First scale down the whole cube by $12$. Let point $P$ have coordinates $(x, y, z)$, point $A$ have coordinates $(0, 0, 0)$, and $s$ be the side length. Then we have the equations
 \begin{align*} (s-x)^2+y^2+z^2&=\left(5\sqrt{10}\right)^2, \\ x^2+(s-y)^2+z^2&=\left(5\sqrt{5}\right)^2, \\ x^2+y^2+(s-z)^2&=\left(10\sqrt{2}\right)^2, \\ (s-x)^2+(s-y)^2+(s-z)^2&=\left(3\sqrt{7}\right)^2. \end{align*}
These simplify into
 \begin{align*} s^2+x^2+y^2+z^2-2sx&=250, \\ s^2+x^2+y^2+z^2-2sy&=125, \\ s^2+x^2+y^2+z^2-2sz&=200, \\ 3s^2-2s(x+y+z)+x^2+y^2+z^2&=63. \end{align*}
Adding the first three equations together, we get $3s^2-2s(x+y+z)+3(x^2+y^2+z^2)=575$.Subtracting this from the fourth equation, we get $2(x^2+y^2+z^2)=512$, so $x^2+y^2+z^2=256$. This means $PA=16$. However, we scaled down everything by $12$ so our answer is $16*12=\boxed{192}$.

 





\textbf{Solution 2 (Solution 1 with Slight Simplification)}

Once the equations for the distance between point P and the vertices of the cube have been written, we can add the first, second, and third to receive, 
 \[2(x^2 + y^2 + z^2) + (s-x)^2 + (s-y)^2 + (s-z)^2 = 250 + 125 + 200.\] Subtracting the fourth equation gives
 \begin{align*} 2(x^2 + y^2 + z^2) &= 575 - 63 \\ x^2 + y^2 + z^2 &= 256 \\ \sqrt{x^2 + y^2 + z^2} &= 16. \end{align*}
Since point $A = (0,0,0), PA = 16$, and since we scaled the answer is $16 \cdot 12 = \boxed{192}$.

 





\textbf{Solution 3}

Let $E$ be the vertex of the cube such that $ABED$ is a square.Using the \href{https://artofproblemsolving.com/wiki/index.php?title=British\_Flag\_Theorem}{British Flag Theorem}, we can easily show that
 \[PA^2 + PE^2 = PB^2 + PD^2\]and
 \[PA^2 + PG^2 = PC^2 + PE^2\]Hence, by adding the two equations together, we get $2PA^2 + PG^2 = PB^2 + PC^2 + PD^2$. Substituting in the values we know, we get $2PA^2 + 7\cdot 36^2 =10\cdot60^2 + 5\cdot 60^2 + 2\cdot 120^2$.

 Thus, we can solve for $PA$, which ends up being $\boxed{192}$.

 (Lokman GÖKÇE)
                    \begin{figure}[H]
                        \centering
                        \includegraphics[width=0.5\textwidth]{Solution/AIME_2021_I_6_Solution_Figure_1.jpg}
                    \end{figure}
                    

 





\textbf{Solution 4}

For all points $X$ in space, define the function $f:\mathbb{R}^{3}\rightarrow\mathbb{R}$ by $f(X)=PX^{2}-GX^{2}$. Then $f$ is linear; let $O=\frac{2A+G}{3}$ be the center of $\triangle BCD$. Then since $f$ is linear,
 \begin{align*} 3f(O)=f(B)+f(C)+f(D)&=2f(A)+f(G) \\  \left(PB^{2}-GB^{2}\right)+\left(PC^{2}-GC^{2}\right)+\left(PD^{2}-GD^{2}\right)&=2\left(PA^{2}-GA^{2}\right)+PG^{2} \\ \left(60\sqrt{10}\right)^{2}-2x^{2}+\left(60\sqrt{5}\right)^{2}-2x^{2}+\left(120\sqrt{2}\right)^{2}-2x^{2}&=2PA^{2}-2\cdot 3x^{2}+\left(36\sqrt{7}\right)^{2}, \end{align*}
where $x$ denotes the side length of the cube. Thus
 \begin{align*}  36\text{,}000+18\text{,}000+28\text{,}800-6x^{2}&=2PA^{2}-6x^{2}+9072 \\  82\text{,}800-6x^{2}&=2PA^{2}-6x^{2}+9072 \\  73\text{,}728&=2PA^{2} \\  36\text{,}864&=PA^{2} \\  PA&=\boxed{192}. \end{align*}


 


 




	\item (\textit{AIME I 2021}) Find the number of pairs $(m,n)$ of positive integers with $1\le m<n\le 30$ such that there exists a real number $x$ satisfying
 \[\sin(mx)+\sin(nx)=2.\]


 
 \textbf{Answer: }063



\textbf{Solution 1}

It is trivial that the maximum value of $\sin \theta$ is $1$, is achieved at $\theta = \frac{\pi}{2}+2k\pi$ for some integer $k$.

 This implies that $\sin(mx) = \sin(nx) = 1$, and that $mx = \frac{\pi}{2}+2a\pi$ and $nx = \frac{\pi}{2}+2b\pi$, for integers $a, b$.

 Taking their ratio, we have 
 \[\frac{mx}{nx} = \frac{\frac{\pi}{2}+2a\pi}{\frac{\pi}{2}+2b\pi} \implies \frac{m}{n} = \frac{4a + 1}{4b + 1} \implies \frac{m}{4a + 1} = \frac{n}{4b + 1} = k.\]It remains to find all $m, n$ that satisfy this equation.

 If $k = 1$, then $m \equiv n \equiv 1 \pmod 4$. This corresponds to choosing two elements from the set $\{1, 5, 9, 13, 17, 21, 25, 29\}$. There are $\binom 82$ ways to do so.

 If $k < 1$, by multiplying $m$ and $n$ by the same constant $c = \frac{1}{k}$, we have that $mc \equiv nc \equiv 1 \pmod 4$. Then either $m \equiv n \equiv 1 \pmod 4$, or $m \equiv n \equiv 3 \pmod 4$. But the first case was already counted, so we don't need to consider that case. The other case corresponds to choosing two numbers from the set $\{3, 7, 11, 15, 19, 23, 27\}$. There are $\binom 72$ ways here. (This argument seems to have a logical flaw *check note at bottom*)

 Finally, if $k > 1$, note that $k$ must be an integer. This means that $m, n$ belong to the set $\{k, 5k, 9k, \dots\}$, or $\{3k, 7k, 11k, \dots\}$. Taking casework on $k$, we get the sets $\{2, 10, 18, 26\}, \{6, 14, 22, 30\}, \{4, 20\}, \{12, 28\}$. Some sets have been omitted; this is because they were counted in the other cases already. This sums to $\binom 42 + \binom 42 + \binom 22 + \binom 22$.

 In total, there are $\binom 82 + \binom 72 + \binom 42 + \binom 42 + \binom 22 + \binom 22 = \boxed{063}$ pairs of $(m, n)$.

 This solution was brought to you by ~Leonard\_my\_dude~

 


 \textbf{Note: Detailed Explanation for $m \equiv n \equiv 3 \pmod 4$} Many of the solutions I see here lack an explicit explanation for why every solution to $m \equiv n \equiv 3 \pmod 4$ constitutes a valid solution to the problem, so I will explain it here. The reason why all solutions where $m \equiv n \equiv 1 \pmod 4$ works is pretty intuitive, as $\frac{m}{n}$ would then be able to be expressed as $\frac{4a + 1}{4b + 1}$. Similarly, note that $-(4x + 1) \equiv 3 \pmod 4$. Thus, when solving the equation for $\frac{m}{n} = \frac{4a + 1}{4b + 1}$, every case in which $m \equiv n \equiv 3 \pmod 4$ will produce a valid solution when $a$ and $b$ are both negative. For example, when $m = 7$ and $n = 11$, $a = -2$ and $b = -3$ gives $\frac{4a + 1}{4b + 1} = \frac{-7}{-11} = \frac{7}{11}$. Thus, we add $\binom 72$ as every case where $m \equiv n \equiv 3 \pmod 4$ produces a valid solution.

 - yangomango

 


 \textbf{Other explanation} Clearly, $mx = 90+360k$ and $nx=90+360a$. From here, we can divide both equations to get $\frac{n}{m} = \frac{1+4a}{1+4k}.$ But $m$ and $n$ might share a common factor, so let that common factor be $d.$ We get $dn \equiv dm \equiv 1 (\bmod 4).$ Taking casework on $d \bmod 4,$ we see either $d \equiv 1 (\bmod 4)$ or $d \equiv 3 (\bmod 4).$ Substituting back, we see $3n \equiv 3m \equiv 1 (\bmod 4)$ or $n \equiv m \equiv 1 (\bmod 4).$ The first case gives $n \equiv m \equiv 3 (\bmod 4),$ and the second obviously gives $n \equiv m \equiv 1 (\bmod 4).$ Now you may proceed as other solutions do. 

 





\textbf{Solution 2}

In order for $\sin(mx) + \sin(nx) = 2$, $\sin(mx) = \sin(nx) = 1$. 

 This happens when $mx \equiv nx \equiv \frac{\pi}{2} ($mod $2\pi).$

 This means that $mx = \frac{\pi}{2} + 2\pi\alpha$ and $nx = \frac{\pi}{2} + 2\pi\beta$ for any integers $\alpha$ and $\beta$. 

 As in Solution 1, take the ratio of the two equations:
 \[\frac{mx}{nx} = \frac{\frac{\pi}{2}+2\pi\alpha}{\frac{\pi}{2}+2\pi\beta} \implies \frac{m}{n} = \frac{\frac{1}{2}+2\alpha}{\frac{1}{2}+2\beta} \implies \frac{m}{n} = \frac{4\alpha+1}{4\beta+1}\]

 Now notice that the numerator and denominator of $\frac{4\alpha+1}{4\beta+1}$ are both odd, which means that $m$ and $n$ have the same power of two (the powers of 2 cancel out).

 Let the common power be $p$: then $m = 2^p\cdot a$, and $n = 2^p\cdot b$ where $a$ and $b$ are integers between 1 and 30.

 We can now rewrite the equation: 
 \[\frac{2^p\cdot a}{2^p\cdot b} = \frac{4\alpha+1}{4\beta+1} \implies \frac{a}{b} = \frac{4\alpha+1}{4\beta+1}\]

 Now it is easy to tell that $a \equiv 1 ($mod $4)$ and $b \equiv 1 ($mod $4)$. However, there is another case: that 

 $a \equiv 3 ($mod $4)$ and $b \equiv 3 ($mod $4)$. This is because multiplying both $4\alpha+1$ and $4\beta+1$ by $-1$ will not change the fraction, but each congruence will be changed to $-1 ($mod $4) \equiv 3 ($mod $4)$.

 From the first set of congruences, we find that $a$ and $b$ can be two of $\{1, 5, 9, \ldots, 29\}$.

 From the second set of congruences, we find that $a$ and $b$ can be two of $\{3, 7, 11, \ldots, 27\}$.

 Now all we have to do is multiply by $2^p$ to get back to $m$ and $n$. Let’s organize the solutions in order of increasing values of $p$, keeping in mind that $m$ and $n$ are bounded between 1 and 30.

 For $p = 0$ we get $\{1, 5, 9, \ldots, 29\}, \{3, 7, 11, \ldots, 27\}$.

 For $p = 1$ we get $\{2, 10, 18, 26\}, \{6, 14, 22, 30\}$

 For $p = 2$ we get $\{4, 20\}, \{12, 28\}$

 Note that $16\mid{a}$ since $m$ will cancel out a factor of 4 from $a$, and $\frac{a}{m}$ must contain a factor of 4. Again, $1-4X$ will never contribute a factor of 2. Simply inspecting, we see two feasible values for $a$ and $m$ such that $a+m\leq30$. 

 If we increase the value of $p$ more, there will be less than two integers in our sets, so we are done there.

 There are 8 numbers in the first set, 7 in the second, 4 in the third, 4 in the fourth, 2 in the fifth, and 2 in the sixth.

 In each of these sets we can choose 2 numbers to be $m$ and $n$ and then assign them in increasing order. Thus there are:

 
 \[\dbinom{8}{2}+\dbinom{7}{2}+\dbinom{4}{2}+\dbinom{4}{2}+\dbinom{2}{2}+\dbinom{2}{2} = 28+21+6+6+1+1 = \boxed{063}\] possible pairs $(m,n)$ that satisfy the conditions.

 -KingRavi

 





\textbf{Solution 3}

We know that the range of sine is between $-1$ and $1$, inclusive.

 Thus, the only way for the sum to be $2$ is for $\sin(mx)=\sin(nx)=1$.

 Note that $\sin(90+360k)=1$.

 Assuming $mx$ and $nx$ are both positive, $m$ and $n$ could be $1,5,9,13,17,21,25,29$. There are $8$ ways, so $\dbinom{8}{2}$.

 If both are negative, $m$ and $n$ could be $3,7,11,15,19,23,27$. There are $7$ ways, so $\dbinom{7}{2}$.

 However, the pair $(1,5)$ could also be $(2, 10)$ and so on. The same goes for some other pairs.

 In total there are $14$ of these extra pairs.

 The answer is $28+21+14 = \boxed{063}$.

 





\textbf{Solution 4}

The equation implies that $\sin(mx)=\sin(nx)=1$. Therefore, we can write $mx$ as $2{\pi}k_1+\frac{\pi}{2}$ and $nx$ as $2{\pi}k_2+\frac{\pi}{2}$ for integers $k_1$ and $k_2$. Then, $\frac{mx}{nx}=\frac{m}{n}=\frac{2k_1+\frac{1}{2}}{2k_2+\frac{1}{2}}$. Cross multiplying, we get $m\cdot{(2k_2+\frac{1}{2})}=n\cdot{(2k_1+\frac{1}{2})} \Longrightarrow 4k_2m-4k_1n=n-m$. Let $n-m=a$ so the equation becomes $4(m(k_2-k_1)+k_1a)=a$. Let $k_2-k_1=X$ and $k_1=Y$, then the equation becomes $a=4Ym+4Xa \Longrightarrow \frac{a(1-4X)}{m}=4Y$. Note that $X$ and $Y$ can vary accordingly, and $4\mid{a}$. Next, we do casework on $m\pmod{4}$:

 If $m\equiv 1\pmod{4}$:

 Once $a$ and $m$ are determined, $n$ is determined, so $a+m\leq30$. $a\in \{4,8,12,\dots,28\}$ and $m\in \{1,5,9,\dots,29\}$. Therefore, there are $\sum_{i=1}^{7}{i}=28$ ways for this case such that $a+m\leq30$.

 If $m\equiv 3\pmod{4}$:

 $a\in \{4,8,12,\dots,28\}$ and $m\in \{3,7,11,\dots,27\}$. Therefore, there are $\sum_{i=1}^{6}{i}=21$ ways such that $a+m\leq30$.

 If $m\equiv 2\pmod{4}$:

 Note that $8\mid{a}$ since $m$ in this case will have a factor of 2, which will cancel out a factor of 2 in $a$, and we need the left hand side to divide 4. Also, $1-4X\equiv 1\pmod{4}$ so it is odd and will therefore never contribute a factor of 2. $a\in \{8,16,24\}$ and $m\in \{2,6,10,\dots,30\}$. Following the condition $a+m\leq30$, we conclude that there are $6+4+2=12$ ways for this case.

 If $m\equiv 0\pmod{4}$:

 Adding all the cases up, we obtain $28+21+12+2=\boxed{063}$

 





\textbf{Solution 5}

It is obvious that $f(x) = \sin(a) + \sin(b) \leq 2$ when equality holds if $a = \frac{\pi}{2} + 2\pi k_1$ and $b = \frac{\pi}{2} + 2\pi k_2$. Therefore, we have

 $mx = \frac{\pi}{2} + 2\pi k_1 \implies x = \frac{\frac{\pi}{2} + 2\pi k_1}{m}$. Similarly, $x = \frac{\frac{\pi}{2} + 2\pi k_2}{n}$. Thus, we can set them equal and cancel to get

 $\frac{\frac{1}{2} + 2k_1}{m} = \frac{\frac{1}{2} + 2k_2}{n} \implies \frac{m}{n} = \frac{4k_1 + 1}{4k_2 + 1}$. We consider $m = 1$ which is already in the form $4k_1 + 1$. Thus, since $m < n$, we just need $n$ to be all the numbers that are in $1 \mod 4$. We have

 $m = 1, n = 5, 9, 13, 17, 21, 25, 29$

 Now, consider $m = 2$. Well, the only way we can make $2$ a number that is $1 \mod 4$ is if we divide $m$ by $2$. This means all the numbers we found earlier for $m = 1$ needs to be multiplied by $2$ to get $n$ into a number that is $1 \mod 4$ when $\frac{m}{n}$ is reduced to lowest terms(Essentially, we're saying $\gcd(m, n) = 2)$). But also, $m < n \leq 30$ so the new numbers we have are

 $m = 2, n = 10, 18, 26$

 For $m = 3$, we find that we just need to multiply $n$ by $3$ and we get

 $m = 3, n = 15, 27$

 For $m = 4$, we multiply by $4$ to get

 $m = 4, n = 20$

 Thus, for $m = 1, 2, 3, 4$ the total number of pairs is $7 + 3 + 2 + 1 = 13$.

 Now, from $m = 5, 6, 7, 8$ note that these residues modulo $4$ are exactly the same as $m = 1, 2, 3, 4$. So we assume a cycle with the same number of ordered pairs for $n$. However, this is not true as some of the cases will lead $m = n$ or $m > n$, We will deal with this later. For now, let's assume we have $13$ ways for $n$ for every $m$ in modulo $4$. We have

 $1 - 4$$5 - 8$$9 - 12$$13 - 16$$17 - 20$$21 - 24$$25 - 28$$29, 30$

 The first $7$ cases give $13$ ways for $n$. The next $2$ cases give us $7 + 3 = 10$ cases considering for every $m$ that is $1 \mod 4$, we have $7$ cases and vice versa. Thus, we have a total of $13 \cdot 7 + 10 = 101$ cases. Now, we need to account for the cases where $m = n$ or $m > n$. We're fine from $m = 1$ to $m = 4$. From $m = 5$ to $m = 8$, we can make a chart.

 $m \equiv 1 \mod 4 \implies n = 5, 9, 13, 17, 21, 25, 29$$m \equiv 2 \mod 4 \implies n = 10, 18, 26$$m \equiv 3 \mod 4 \implies n = 15, 27$$m \equiv 0 \mod 4 \implies n = 20$

 For $m = 5$, we know this is the case for $m \equiv 1 \mod 4$. So we can clearly see that the case where $n =  5$ doesn't work. Thus $(m, n) = (5, 5)$ doesn't work. For $m = 6, m = 7, m = 8$ we can look at the chart and see everything else works.

 For $m = 9$, we see $n = 5, n = 9$ doesn't work. For $m = 10$ we see $n = 10$ doesn't work.

 Following this fashion we get the following pairs don't work

 $(m, n) = (5, 5), (9, 5), (9, 9), (10, 10), (13, 5), (13, 9), (13, 13), (14, 10), (15, 15), (17, 5), (17, 9), (17, 13), (17, 17), (18, 10), (18, 18), (19, 15), (20, 20)$ and $(21, 5), (21, 9), (21, 13), (21, 17), (21, 21)$ and $(22, 10), (22, 18), (23, 15), (24, 20), (25, 5), (25, 9), (25, 13), (25, 17), (25, 21), (25, 25)$ and $(26, 10) (26, 18), (26, 20), (27, 15), (27, 27), (28, 20)$. This gives $38$ cases when counting.

 Thus, the answer is $101 - 38 = \boxed{063}$

 \textbf{Remark} The graphs of $r\leq\sin(m\theta)+\sin(n\theta)$ and $r=2$ are shown here in Desmos: \href{https://www.desmos.com/calculator/busxadywja}{https://www.desmos.com/calculator/busxadywja}

 Move the sliders around for $1\leq m \leq 29$ and $2\leq m+1\leq n\leq30$ to observe the geometric representation generated by each pair $(m,n).$

 




	\item (\textit{AIME I 2021}) Find the number of integers $c$ such that the equation
 \[\left||20|x|-x^2|-c\right|=21\]has $12$ distinct real solutions.


 
 \textbf{Answer: }057



\textbf{Solution 1 (Piecewise Function: Analysis and Graph)}

We take cases for the outermost absolute value, then rearrange: 
 \[\left|20|x|-x^2\right|=c\pm21.\]Let $f(x)=\left|20|x|-x^2\right|.$ We rewrite $f(x)$ as a piecewise function without using absolute values:
 \[f(x) = \begin{cases} \left|-20x-x^2\right| & \mathrm{if} \ x \le 0  \begin{cases} 20x+x^2 & \mathrm{if} \ x\le-20 \\ -20x-x^2 & \mathrm{if} \ -20<x\leq0 \end{cases} \\ \left|20x-x^2\right| & \mathrm{if} \ x > 0 \begin{cases} 20x-x^2 & \mathrm{if} \ 0<x\leq20 \\ -20x+x^2 & \mathrm{if} \ x>20 \end{cases} \end{cases}.\]We graph $y=f(x)$ with all extremum points labeled, as shown below. The fact that $f(x)$ is an even function ($f(x)=f(-x)$ holds for all real numbers $x,$ so the graph of $y=f(x)$ is symmetric about the $y$-axis) should facilitate the process of graphing.
                    \begin{figure}[H]
                        \centering
                        \includegraphics[width=0.5\textwidth]{Solution/AIME_2021_I_8_Solution_Figure_1.png}
                    \end{figure}
                    Since $f(x)=c\pm21$ has $12$ distinct real solutions, it is clear that each case has $6$ distinct real solutions geometrically. We shift the graphs of $y=\pm21$ up $c$ units, where $c\geq0:$

 
                \begin{itemize}
                    \item For $f(x)=c+21$ to have $6$ distinct real solutions, we need $0\leq c<79.$
                \end{itemize}


                
                \begin{itemize}
                    \item For $f(x)=c-21$ to have $6$ distinct real solutions, we need $21<c<121.$
                \end{itemize}


                Taking the intersection of these two cases gives $21<c<79,$ from which there are $79-21-1=\boxed{057}$ such integers $c.$

 





\textbf{Solution 2 (Graphing)}

Graph $y=|20|x|-x^2|$ (If you are having trouble, look at the description in the next two lines and/or the diagram in Solution 1). Notice that we want this to be equal to $c-21$ and $c+21$.

 We see that from left to right, the graph first dips from very positive to $0$ at $x=-20$, then rebounds up to $100$ at $x=-10$, then falls back down to $0$ at $x=0$.

 The positive $x$ are symmetric, so the graph re-ascends to $100$ at $x=10$, falls back to $0$ at $x=10$, and rises to arbitrarily large values afterwards.

 Now we analyze the $y$ (varied by $c$) values. At $y=k<0$, we will have no solutions, as the line $y=k$ will have no intersections with our graph.

 At $y=0$, we will have exactly $3$ solutions for the three zeroes.

 At $y=n$ for any $n$ strictly between $0$ and $100$, we will have exactly $6$ solutions.

 At $y=100$, we will have $4$ solutions, because local maxima are reached at $x= \pm 10$.

 At $y=m>100$, we will have exactly $2$ solutions.

 To get $12$ distinct solutions for $y=|20|x|-x^2|=c \pm 21$, both $c +21$ and $c-21$ must produce $6$ solutions.

 Thus $0<c-21$ and $c+21<100$, so $c \in \{ 22, 23, \dots , 77, 78 \}$ is required.

 It is easy to verify that all of these choices of $c$ produce $12$ distinct solutions (none overlap), so our answer is $\boxed{057}$.

 





\textbf{Solution 3 (Graphing)}

Let $y = |x|.$ Then the equation becomes $\left|\left|20y-y^2\right|-c\right| = 21$, or $\left|y^2-20y\right| = c \pm 21$. Note that since $y = |x|$, $y$ is nonnegative, so we only care about nonnegative solutions in $y$. Notice that each positive solution in $y$ gives two solutions in $x$ ($x = \pm y$), whereas if $y = 0$ is a solution, this only gives one solution in $x$, $x = 0$. Since the total number of solutions in $x$ is even, $y = 0$ must not be a solution. Hence, we require that $\left|y^2-20y\right| = c \pm 21$ has exactly $6$ positive solutions and is not solved by $y = 0.$

 If $c < 21$, then $c - 21$ is negative, and therefore cannot be the absolute value of $y^2 - 20y$. This means the equation's only solutions are in $\left|y^2-20y\right| = c + 21$. There is no way for this equation to have $6$ solutions, since the quadratic $y^2-20y$ can only take on each of the two values $\pm(c + 21)$ at most twice, yielding at most $4$ solutions. Hence, $c \ge 21$. $c$ also can't equal $21$, since this would mean $y = 0$ would solve the equation. Hence, $c > 21.$

 At this point, the equation $y^2-20y = c \pm 21$ will always have exactly $2$ positive solutions, since $y^2-20y$ takes on each positive value exactly once when $y$ is restricted to positive values (graph it to see this), and $c \pm 21$ are both positive. Therefore, we just need $y^2-20y = -(c \pm 21)$ to have the remaining $4$ solutions exactly. This means the horizontal lines at $-(c \pm 21)$ each intersect the parabola $y^2 - 20y$ in two places. This occurs when the two lines are above the parabola's vertex $(10,-100)$. Hence we have
 \begin{align*} -(c + 21) &> -100 \\ c + 21 &< 100 \\ c &< 79. \end{align*}
Hence, the integers $c$ satisfying the conditions are those satisfying $21 < c < 79.$ There are $\boxed{057}$ such integers.

 Note: Be careful of counting at the end, you may mess up and get $59$.

 





\textbf{Solution 4 (Algebra)}

Removing the absolute value bars from the equation successively, we get
 \begin{align*} \left|\left|20|x|-x^2\right|-c\right|&=21 \\ \left|20|x|-x^2\right|&= c \pm21 \\ 20|x|-x^2 &= \pm c \pm 21 \\ x^2 \pm 20x \pm c \pm21 &= 0. \end{align*}
The discriminant of this equation is
 \[\sqrt{400-4(\pm c \pm 21)}.\]Equating the discriminant to $0$, we see that there will be two distinct solutions to each of the possible quadratics above only in the interval $-79 < c < 79$. However, the number of zeros the equation $ax^2+b|x|+k$ has is determined by where $ax^2+bx+k$ and $ax^2-bx+k$ intersect, namely at $(0,k)$. When $k>0$, $a>0$, $ax^2+b|x|+k$ will have only $2$ solutions, and when $k<0$, $a>0$, then there will be $4$ real solutions, if they exist at all.In order to have $12$ solutions here, we thus need to ensure $-c+21<0$, so that exactly $2$ out of the $4$ possible equations of the form $ax^2+b|x|+k$ given above have y-intercepts below $0$ and only $2$ real solutions, while the remaining $2$ equations have $4$ solutions. This occurs when $c>21$, so our final bounds are $21<c<79$, giving us $\boxed{057}$ valid values of $c$.

 \textbf{Remark} The graphs of $F(x)=\left||20|x|-x^2|-c\right|$ and $G(x)=21$ are shown here in Desmos: \href{https://www.desmos.com/calculator/i6l98lxwpp}{https://www.desmos.com/calculator/i6l98lxwpp}

 Move the slider around for $21<c<79$ to observe how they intersect for $12$ times.

 




	\item (\textit{AIME I 2021}) Let $ABCD$ be an isosceles trapezoid with $AD=BC$ and $AB<CD.$ Suppose that the distances from $A$ to the lines $BC,CD,$ and $BD$ are $15,18,$ and $10,$ respectively. Let $K$ be the area of $ABCD.$ Find $\sqrt2 \cdot K.$


 
 \textbf{Answer: }567



\textbf{Solution 1 (Similar Triangles and Pythagorean Theorem)}

Let $\overline{AE}, \overline{AF},$ and $\overline{AG}$ be the perpendiculars from $A$ to $\overleftrightarrow{BC}, \overleftrightarrow{CD},$ and $\overleftrightarrow{BD},$ respectively. Next, let $H$ be the intersection of $\overline{AF}$ and $\overline{BD}.$

 We set $AB=x$ and $AH=y,$ as shown below.
                    \begin{figure}[H]
                        \centering
                        \includegraphics[width=0.5\textwidth]{Solution/AIME_2021_I_9_Solution_Figure_1.png}
                    \end{figure}
                    From here, we obtain $HF=18-y$ by segment subtraction, and $BG=\sqrt{x^2-10^2}$ and $HG=\sqrt{y^2-10^2}$ by the Pythagorean Theorem.

 Since $\angle ABG$ and $\angle HAG$ are both complementary to $\angle AHB,$ we have $\angle ABG = \angle HAG,$ from which $\triangle ABG \sim \triangle HAG$ by AA. It follows that $\frac{BG}{AG}=\frac{AG}{HG},$ so $BG\cdot HG=AG^2,$ or 
 \[\sqrt{x^2-10^2}\cdot\sqrt{y^2-10^2}=10^2. \hspace{10mm} (1)\]Since $\angle AHB = \angle FHD$ by vertical angles, we have $\triangle AHB \sim \triangle FHD$ by AA, with the ratio of similitude $\frac{AH}{FH}=\frac{BA}{DF}.$ It follows that $DF=BA\cdot\frac{FH}{AH}=x\cdot\frac{18-y}{y}.$

 Since $\angle EBA = \angle ECD = \angle FDA$ by angle chasing, we have $\triangle EBA \sim \triangle FDA$ by AA, with the ratio of similitude $\frac{EA}{FA}=\frac{BA}{DA}.$ It follows that $DA=BA\cdot\frac{FA}{EA}=x\cdot\frac{18}{15}=\frac{6}{5}x.$

 By the Pythagorean Theorem on right $\triangle ADF,$ we have $DF^2+AF^2=AD^2,$ or 
 \[\left(x\cdot\frac{18-y}{y}\right)^2+18^2=\left(\frac{6}{5}x\right)^2. \hspace{7mm} (2)\]Solving this system of equations ($(1)$ and $(2)$), we get $x=\frac{45\sqrt2}{4}$ and $y=\frac{90}{7},$ so $AB=x=\frac{45\sqrt2}{4}$ and $CD=AB+2DF=x+2\left(x\cdot\frac{18-y}{y}\right)=\frac{81\sqrt2}{4}.$ Finally, the area of $ABCD$ is 
 \[K=\frac{AB+CD}{2}\cdot AF=\frac{567\sqrt2}{2},\] from which $\sqrt2 \cdot K=\boxed{567}.$

 <u>\textbf{Remark}</u>

 Instead of solving the system of equations $(1)$ and $(2),$ which can be time consuming, by noting that $\triangle ACF \sim \triangle ABG$ by AA, we could find out $\frac{AB}{AG} = \frac{AC}{AF}$, which gives $AC = \frac{9}{5}x$. We also know that $EB = \sqrt{x^2 - 15^2}$ by Pythagorean Theorem on $\triangle ABE$. From $BC = AD = \frac{6}{5}x,$ we apply the Pythagorean Theorem to $\triangle ACE$ and obtain
 \[AC^2 = (EB+BC)^2 + AE^2.\]Substituting, we get
 \[\frac{81}{25}x^2 = \left(\sqrt{x^2 -225}+\frac{6}{5}x\right)^2+225 \iff x = 3\sqrt{x^2 - 15^2},\]from which $x = \frac{45\sqrt{2}}{4}.$

 





\textbf{Solution 2 (Similar Triangles and Pythagorean Theorem)}

First, draw the diagram. Then, notice that since $ABCD$ is isosceles, $\Delta ABD \cong \Delta BAC$, and the length of the altitude from $B$ to $AC$ is also $10$. Let the foot of this altitude be $F$, and let the foot of the altitude from $A$ to $BC$ be denoted as $E$. Then, $\Delta BCF \sim \Delta ACE$. So, $\frac{BC}{AC} = \frac{BF}{AE} = \frac{2}{3}$. Now, notice that $[ABC] = \frac{10 \cdot AC} {2} = \frac{AB \cdot 18}{2} \implies AC = \frac{9 \cdot AB}{5}$, where $[ABC]$ denotes the area of triangle $ABC$. Letting $AB = x$, this equality becomes $AC = \frac{9x}{5}$. Also, from $\frac{BC}{AC} = \frac{2}{3}$, we have $BC = \frac{6x}{5}$. Now, by the Pythagorean theorem on triangles $ABF$ and $CBF$, we have $AF = \sqrt{x^{2}-100}$ and $CF = \sqrt{ \left( \frac{6x}{5} \right) ^{2}-100}$. Notice that $AC = AF + CF$, so $\frac{9x}{5} = \sqrt{x^{2}-100} + \sqrt{ \left( \frac{6x}{5} \right) ^{2}-100}$. Squaring both sides of the equation once, moving $x^{2}-100$ and $\left( \frac{6x}{5} \right) ^{2}-100$ to the right, dividing both sides by $2$, and squaring the equation once more, we are left with $\frac{32x^{4}}{25} = 324x^{2}$. Dividing both sides by $x^{2}$ (since we know $x$ is positive), we are left with $\frac{32x^{2}}{25} = 324$. Solving for $x$ gives us $x = \frac{45}{2\sqrt{2}}$. 

 Now, let the foot of the perpendicular from $A$ to $CD$ be $G$. Then let $DG = y$. Let the foot of the perpendicular from $B$ to $CD$ be $H$. Then, $CH$ is also equal to $y$. Notice that $ABHG$ is a rectangle, so $GH = x$. Now, we have $CG = GH + CH = x + y$. By the Pythagorean theorem applied to $\Delta AGC$, we have $(x+y)^{2}+18^{2}= \left( \frac{9x}{5} \right) ^{2}$. We know that $\frac{9x}{5} = \frac{9}{5} \cdot \frac{45}{2\sqrt{2}} = \frac{81}{2\sqrt{2}}$, so we can plug this into this equation. Solving for $x+y$, we get $x+y=\frac{63}{2\sqrt{2}}$. 

 Finally, to find $[ABCD]$, we use the formula for the area of a trapezoid: $K = [ABCD] = \frac{b_{1}+b_{2}}{2} \cdot h = \frac{AB+CD}{2} \cdot 18 = \frac{x+(CG+DG)}{2} \cdot 18 = \frac{2x+2y}{2} \cdot 18 = (x+y) \cdot 18 = \frac{63}{2\sqrt{2}} \cdot 18 = \frac{567}{\sqrt{2}}$. The problem asks us for $K \cdot \sqrt{2}$, which comes out to be $\boxed{567}$.

 





\textbf{Solution 3 (Similar Triangles and Pythagorean Theorem)}

Make $AE$ perpendicular to $BC$; $AG$ perpendicular to $BD$; $AF$ perpendicular $DC$.

 It's obvious that $\triangle{AEB} \sim \triangle{AFD}$. Let $EB=5x; AB=5y; DF=6x; AD=6y$. Then make $BQ$ perpendicular to $DC$, it's easy to get $BQ=18$.

 Since $AB$ parallel to $DC$, $\angle{ABG}=\angle{BDQ}$, so $\triangle{ABG} \sim \triangle{BDQ}$. After drawing the altitude, it's obvious that $FQ=AB=5y$, so $DQ=5y+6x$. According to the property of similar triangles, $AG/BQ=BG/DQ$. So, $\frac{5}{9}=\frac{GB}{(6x+5y)}$, or $GB=\frac{(30x+25y)}{9}$.

 Now, we see the $\triangle AEB$, pretty easy to find that $15^2+(5x)^2=(5y)^2$, then we get $x^2+9=y^2$, then express $y$ into $x$ form that $y=\sqrt{x^2+9}$we put the length of $BG$ back to $\triangle AGB$: $BG^2+100=AB^2$. So, 
 \[\frac{[30x+25\sqrt{(x^2+9)}]^2}{81}+100=(5\sqrt{x^2+9})^2.\]After calculating, we can have a final equation of $x^2+9=\sqrt{x^2+9}\cdot3x$. It's easy to find $x=\frac{3\sqrt{2}}{4}$ then $y=\frac{9\sqrt{2}}{4}$. So, 
 \[\sqrt{2}\cdot K = \sqrt{2}\cdot(5y+5y+6x+6x)\cdot9=\boxed{567}.\]~bluesoul

 





\textbf{Solution 4 (Cyclic Quadrilaterals, Similar Triangles, Pythagorean Theorem)}

Let the foot of the altitude from $A$ to $BC$ be $P$, to $CD$ be $Q$, and to $BD$ be $R$.

 Note that all isosceles trapezoids are cyclic quadrilaterals; thus, $A$ is on the circumcircle of $\triangle BCD$ and we have that $PRQ$ is the Simson Line from $A$. As $\angle QAB = 90^\circ$, we have that $\angle QAR = 90^\circ - \angle RAB =\angle ABR = \angle APR = \angle APQ$, with the last equality coming from cyclic quadrilateral $APBR$. Thus, $\triangle QAR \sim \triangle QPA$ and we have that $\frac{AQ}{AR} = \frac{PQ}{PA}$ or that $\frac{18}{10} = \frac{QP}{15}$, which we can see gives us that $QP = 27$. Further ratios using the same similar triangles gives that $QR = 12$ and $RP = 15$.

 We also see that quadrilaterals $APBR$ and $ARQD$ are both cyclic,it is clear that $\angle AQR =\angle ADR$ and $\angle APR =\angle ABR$,which shows $\triangle ABD \sim \triangle APQ$. we wish to find the ratio of similitude between the two triangles.

 To do this, we use the one number we have for $\triangle ABD$: we know that the altitude from $A$ to $BD$ has length $10$. As the two triangles are similar, if we can find the height from $A$ to $PQ$, we can take the ratio of the two heights as the ratio of similitude. To do this, we once again note that $QP = 27$. Using this, we can drop the altitude from $A$ to $QP$ and let it intersect $QP$ at $H$. Then, let $QH = x$ and thus $HP=27-x$. We then have by the Pythagorean Theorem on $\triangle AQH$ and $\triangle APH$:
 \begin{align*} 18^2 - x^2 &= 15^2 - (27-x)^2 \\ 324 - x^2 &= 225 - (x^2-54x+729) \\ 54x &= 828 \\ x &= \frac{46}{3}. \end{align*}
Then, $RH = QH - QR = \frac{46}{3} - 12 = \frac{10}{3}$. This gives us then from right triangle $\triangle ARH$ that $AH = \frac{20\sqrt{2}}{3}$ and thus the ratio of $\triangle APQ$ to $\triangle ABD$ is $\frac{3\sqrt{2}}{4}$. From this, we see then that 
 \[AB = AP \cdot \frac{3\sqrt{2}}{4} = 15 \cdot \frac{3\sqrt{2}}{4} = \frac{45\sqrt{2}}{4}\] and 
 \[AD = AQ \cdot \frac{3\sqrt{2}}{4} = 18 \cdot \frac{3\sqrt{2}}{4} = \frac{27\sqrt{2}}{2}.\] The Pythagorean Theorem on $\triangle AQD$ then gives that 
 \[QD = \sqrt{AD^2 - AQ^2} = \sqrt{\left(\frac{27\sqrt{2}}{2}\right)^2 - 18^2} = \sqrt{\frac{81}{2}} = \frac{9\sqrt{2}}{2}.\]Then, we have the height of trapezoid $ABCD$ is $AQ = 18$, the top base is $AB = \frac{45\sqrt{2}}{4}$, and the bottom base is $CD = \frac{45\sqrt{2}}{4} + 2\cdot\frac{9\sqrt{2}}{2}$. From the equation of a trapezoid, $K = \frac{b_1+b_2}{2} \cdot h = \frac{63\sqrt{2}}{4} \cdot 18 = \frac{567\sqrt{2}}{2}$, so the answer is $K\sqrt{2} = \boxed{567}$.

 





\textbf{Solution 5 (Cyclic Quadrilaterals, Similar Triangles, Pythagorean Theorem)}

Let $E,F,$ and $G$ be the feet of the altitudes from $A$ to $BC,CD,$ and $DB$, respectively.

 Claim: We have $2$ pairs of similar right triangles: $\triangle AEB \sim \triangle AFD$ and $\triangle AGD \sim \triangle AEC$.

 Proof: Note that $ABCD$ is cyclic. We need one more angle, and we get this from this cyclic quadrilateral: 
 \begin{align*} \angle ABE &= 180^\circ - \angle ABC =\angle ADC = \angle ADG, \\ \angle ADG &= \angle ADB =\angle ACB = \angle ACE. \hspace{20mm} \square \end{align*}
Let $AD=a$. We obtain from the similarities $AB = \frac{5a}{6}$ and $AC=BD=\frac{3a}{2}$.

 By Ptolemy, $\left(\frac{3a}{2}\right)^2 = a^2 + \frac{5a}{6} \cdot CD$, so $\frac{5a^2}{4} = \frac{5a}{6} \cdot CD$.

 We obtain $CD=\frac{3a}{2}$, so $DF=\frac{CD-AB}{2}=\frac{a}{3}$.

 Applying the Pythagorean theorem on $\triangle ADF$, we get $324=a^2 - \frac{a^2}{9}=\frac{8a^2}{9}$.

 Thus, $a=\frac{27}{\sqrt{2}}$, and $[ABCD]=\frac{AB+CD}{2} \cdot 18 = \frac{\frac{5a}{6} +\frac{9a}{6}}{2} \cdot 18 = 18 \cdot \frac{7}{6} \cdot \frac{27}{\sqrt{2}} = \frac{567}{\sqrt{2}}$, yielding $\sqrt2\cdot[ABCD]=\boxed{567}$.

 





\textbf{Solution 6 (Similar Triangles and Trigonometry)}

Let $AD=BC=a$. Draw diagonal $AC$ and let $G$ be the foot of the perpendicular from $B$ to $AC$, $F$ be the foot of the perpendicular from $A$ to line $BC$, and $H$ be the foot of the perpendicular from $A$ to $DC$.

 Note that $\triangle CBG\sim\triangle CAF$, and we get that $\frac{10}{15}=\frac{a}{AC}$. Therefore, $AC=\frac32 a$. It then follows that $\triangle ABF\sim\triangle ADH$. Using similar triangles, we can then find that $AB=\frac{5}{6}a$. Using the Law of Cosines on $\triangle ABC$, We can find that the $\cos\angle ABC=-\frac{1}{3}$. Since $\angle ABF=\angle ADH$, and each is supplementary to $\angle ABC$, we know that the $\cos\angle ADH=\frac{1}{3}$. It then follows that $a=\frac{27\sqrt{2}}{2}$. Then it can be found that the area $K$ is $\frac{567\sqrt{2}}{2}$. Multiplying this by $\sqrt{2}$, the answer is $\boxed{567}$.

 





\textbf{Solution 7 (Similar Triangles and Trigonometry)}

Draw the distances in terms of $B$, as shown in the diagram. By similar triangles, $\triangle{AEC}\sim\triangle{BIC}$. As a result, let $AB=u$, then $BC=AD=\frac{6}{5}u$ and $2AC=3BC$. The triangle $ABC$ is $6-5-9$ which $\cos(\angle{ABC})=-\frac{1}{3}$. By angle subtraction, $\cos(180-\theta)=-\cos\theta$. Therefore, $AB=\frac{45}{2\sqrt{2}}=\frac{45\sqrt{2}}{4}$ and $AD=BC=\frac{27}{\sqrt{2}}$. By trapezoid area formula, the area of $ABCD$ is equal to $(AB+DF)\cdot 18=567\cdot \frac{\sqrt{2}}{2}$ which $\sqrt{2}\cdot k=\boxed{567}$.

 





\textbf{Solution 8 (Heron's Formula)}


                    \begin{figure}[H]
                        \centering
                        \includegraphics[width=0.5\textwidth]{Solution/AIME_2021_I_9_Solution_Figure_2.png}
                    \end{figure}
                    Let the points formed by dropping altitudes from $A$ to the lines $BC$, $CD$, and $BD$ be $E$, $F$, and $G$, respectively.

 We have
 \[\triangle ABE \sim \triangle ADF \implies \frac{AD}{18} = \frac{AB}{15} \implies AD = \frac{6}{5}AB\]and
 \[BD\cdot10 = 2[ABD] = AB\cdot18 \implies BD = \frac{9}{5}AB.\]For convenience, let $AB = 5x$. By Heron's formula on $\triangle ABD$, we have sides $5x,6x,9x$ and semiperimeter $10x$, so
 \[\sqrt{10x\cdot5x\cdot4x\cdot1x} = [ABD] = \frac{AB\cdot18}{2} = 45x \implies 10\sqrt{2}x^2 = 45x \implies x= \frac{45}{10\sqrt{2}},\]so $AB = 5x = \frac{45}{2\sqrt{2}}$.

 Then,
 \[BE = \sqrt{AB^2 - CA^2} = \sqrt{\left(\frac{45}{2\sqrt{2}}\right)^2 - 15^2} = \sqrt{\frac{225}{8}} = \frac{15}{2\sqrt{2}}\]and
 \[\triangle ABE \sim \triangle ADF \implies DF = \frac{6}{5}BE = \frac{6}{5}\cdot\frac{15}{2\sqrt{2}} = \frac{18}{2\sqrt{2}}.\]Finally, recalling that $ABCD$ is isosceles,
 \[K = [ABCD] = \frac{18}{2}(AB + (AB + 2DF)) = 18(AB + DF) = 18\left(\frac{45}{2\sqrt{2}} + \frac{18}{2\sqrt{2}}\right) = \frac{567}{\sqrt{2}},\]so $\sqrt{2}\cdot K = \boxed{567}$.

 





\textbf{Solution 9 (Three Heights)}


                    \begin{figure}[H]
                        \centering
                        \includegraphics[width=0.5\textwidth]{Solution/AIME_2021_I_9_Solution_Figure_3.jpg}
                    \end{figure}
                    Let $\overline{AE}, \overline{AF},$ and $\overline{AG}$ be the perpendiculars from $A$ to ${BC}, {CD},$ and ${BD},$ respectively.$AE = 15, AF = 18, AG =10$.Denote by $G'$ the base of the perpendicular from $B$ to $AC, H$ be the base of the perpendicular from $C$ to $AB$. Denote $\theta = \angle{CBH}.$It is clear that 
 \[BG' = AG, CH = AF, \triangle CBH \ =\triangle ADF,\] the area of $ABCD$  is equal to the area of the rectangle   $AFCH.$

 The problem is reduced to finding $AH$.

 In triangle $ABC$ all altitudes are known: 
 \[AB : BC : AC =  \frac{1}{CH}\ : \frac{1}{AE}\ : \frac{1}{BG'}\ =\]
 \[= \frac{1}{AF}\ :  \frac{1}{AE}\ :  \frac{1}{AG}\ = 5 : 6 : 9.\]We apply the Law of Cosines to $\triangle ABC$ and get$:$
 \begin{align*} 2\cdot AB\cdot BC \cdot \cos\theta = AC^2 – AB^2 – BC^2, \end{align*}

 \begin{align*} 2\cdot 5\cdot 6\cdot \cos\theta  = 60 \cos\theta  = 9^2 – 5^2 – 6^2 = 20, \cos\theta  =\frac{1}{3}. \end{align*}

 \begin{align*} BH = BC \cos\theta = \frac{BC}{3}.\end{align*}
We apply the Pythagorean Law to $\triangle HBC$ and get$:$
 \begin{align*} HC^2 = 18^2 = BC^2 – BH^2 = 9\cdot BH^2 – BH^2 = 8 BH^2.\end{align*}

 \begin{align*} BH =  \frac{9}{\sqrt2}, AH = (\frac{5}{2} + 1)\cdot BH =  \frac{63}{2\cdot \sqrt2}. \end{align*}
  Required area is
 \begin{align*}  K =  \frac{63}{2\cdot \sqrt{2}} \cdot 18 =  \frac{567}{\sqrt{2}} \implies \sqrt{2} K=\boxed{567}. \end{align*}
 

 \textbf{vladimir.shelomovskii@gmail.com, vvsss}

 





\textbf{Solution 10 (Area)}

Let $F$ be on $DC$ such that $AF \| DC$. Let $G$ be on $BD$ such that $AG \| BD$.

 Let $m$ be the length of $AB$. Let $n$ be the length of $AD$.

 The area of $\triangle ABD$ can be expressed in three ways: $\frac{1}{2}(15)(BC) = \frac{1}{2}(15)(n)$, $\frac{1}{2}(18)(m)$, and $\frac{1}{2}(10)(BD)$.

 
 \[\frac{1}{2}(15)(n) = \frac{1}{2}(18)(m)\]
 \[15n = 18m\]
 \[5n = 6m\]
 \[n = \frac{6}{5}m\]

 Now, $BD = BG + GD = \sqrt{m^2-100} + \sqrt{n^2-100}$. We can substitute in $n = \frac{6}{5}m$ to get

 $BD = \sqrt{m^2-100} + \sqrt{(\frac{6}{5}m)^2-100}$. 

 We have
 \[\frac{1}{2}(10)\left(\sqrt{m^2-100} + \sqrt{(\frac{6}{5}m)^2-100}\right) = \frac{1}{2}(18)(m)\]After a fairly straightforward algebraic bash, we get $m = \frac{45\sqrt{2}}{4}$, and $n = (\frac{6}{5})(\frac{45\sqrt{2}}{4}) = \frac{27\sqrt{2}}{2}$. By the Pythagorean Theorem on $\triangle ADF$, $DF^2 = n^2 - 18^2 = \frac{729}{2} - 324 = \frac{81}{2}$, and $DF = \frac{9\sqrt{2}}{2}$.

 Thus, $DC = 2DF + AB = 9\sqrt{2}+\frac{45\sqrt{2}}{4} = \frac{81\sqrt{2}}{4}$. Therefore, $K = \frac{1}{2}(\frac{45\sqrt{2}}{4}+\frac{81\sqrt{2}}{4}) \cdot 18 = \frac{63\sqrt{2}}{2} \cdot 18 = \frac{567\sqrt{2}}{2}$. The requested answer is $K \cdot \sqrt{2} = \boxed{567}$.

 


 





\textbf{Solution 11 (Coordinates)}

By the given information, the height of the trapezoid is $18,$ but to avoid numerical calculations, call it $b.$ Label $A(-a, b)$, $B(a, b)$, $C(c, 0)$, $D(-c, 0),$ with $b>0, 0<a<c.$ We have the line equations $BD: y=\frac{b}{a+c}x + \frac{bc}{a+c}, BC: y=\frac{b}{a-c}x - \frac{bc}{a-c},$ or rewriting them, $BD: bx-(a+c)y + bc=0, BC: bx - (a-c)y-bc = 0.$ Hence applying two projection theorems gives

 
 \[10 = \frac{|-ab - b(a+c) + bc|}{\sqrt{b^2 + (a+c)^2}},\]
 \[15 = \frac{|-ab - b(a-c) - bc|}{\sqrt{b^2 + (a-c)^2}}.\]Simplifying and remembering $a,b>0,$

 
 \[10 = \frac{36a}{\sqrt{324+(a+c)^2}},\]
 \[15 = \frac{36a}{\sqrt{324+(a-c)^2}}.\]

 Hence 
 \[324 + (a+c)^2 = \left(\frac{18a}{5}\right)^2,\]
 \[324 + (a-c)^2 = \left(\frac{12a}{5}\right)^2.\]

 Subtracting the two equations yield $a = \frac{5}{9}c.$ Plugging that back into the first equation gives $324 + \left(\frac{14}{9}c\right)^2 = 4c^2$ or $\frac{128}{81}c^2 = 324, c^2 = \frac{81^2}{32} \rightarrow c = \frac{81}{4\sqrt{2}}.$ Hence $a = \frac{45}{4\sqrt{2}},$ and the area of the trapezoid is $18 \left( \frac{2a + 2c}{2}\right) = 18(a+c) = \frac{18 \cdot 63}{2\sqrt{2}}=\frac{567}{\sqrt{2}}.$ Answer extraction gives $\boxed{567}.$

 


 





\textbf{Solution 12 (Trig Bash Using Heights)}

We will use trig bash. After some trial, we can see that knowing the size of one angle of the trapezoid allows us to calculate every length of it easily. We will try to find $\angle C$ in this solution.

 
Let $\angle C=a$ and $\angle ABD=b$. $\angle ADB=a-b$ as $AB$ and $CD$ are parallel lines.

 
The height of the trapezoid is $18$. By trigs we have $AD=\frac{18}{\sin{a}}$.

 
Inside the triangle $ABD$, its height from point $A$ is $10$, its height from point $A$ is also $AD \sin{\angle ADB}=\frac{18}{\sin{a}}\sin(a-b)$. We obtain the equation $\frac{18}{\sin{a}}\sin(a-b)=10$.

 
By trigs, $AB$ is both $\frac{15}{\sin{a}}$ and $\frac{10}{\sin{b}}$. $\frac{15}{\sin{a}} = \frac{10}{\sin{b}}$.

 
Expanding the equation $\frac{18}{\sin{a}}\sin(a-b)=10$ using angle-sum formula, we have $18(\sin{a}\cos{b}-\cos{a}\sin{b})=10\sin{a}$. We can just plug in $\sin{b}=\frac{2}{3}\sin{a}$ and use the identity $\cos{a}=\sqrt{1-\sin^2{a}}$ now to solve this common, algebraic equation.The process of solving is long and boring. Eventually we will get $\sin{a}=\frac{2\sqrt{2}}{3}$.

 
In the end, by trigs we have$AB=\frac{15}{\sin{a}}$             and             $CD=\frac{15}{\sin{a}}+2\frac{18}{\tan{a}}$

 
The area of $ABCD$ is $(AB+CD)\times \frac{18}{2} = 9(\frac{30}{\sin{a}}+\frac{36}{\tan{a}})$. Plugging in $\sin{a}=\frac{2\sqrt{2}}{3}$ and $\tan{a}=\frac{\sin{a}}{\sqrt{1-\sin^2{a}}}=2\sqrt{2}$, we find that the area of $ABCD$ is $\frac{567\sqrt{2}}{2}$.

 
The answer to the problem is $567$.

 





\textbf{Solution 13}

Let the foot of the altitude from $A$ to $CD$ be $E$. Then $AE = 18$. Let the foot of the altitude from $A$ to $BD$ be $F$. Then $AF = 10$. Finally, let the foot of the altitude from $A$ to $BC$ be $G$. Then $AG = 15$. Also, let the foot of the altitude from $B$ to $CD$ be $H$. Then $AE = BH = 18$. Let $\angle CBH = \theta$. Notice $ABHE$ is a rectangle since $ABCD$ is an isosceles trapezoid. Thus, $\angle HBA = 90^{\circ}$. Notice $B, C, G$ are collinear. Thus, $\angle ABG = 180 - 90 - \theta = 90 - \theta \implies \triangle ABG \sim \triangle BCH$. Note $\triangle BCH \cong \triangle ADE$. Let $DE = CH = t$ and $BG = a$. Then by similar triangles proportions we have

 $\frac{t}{a} = \frac{18}{15} \implies t = \frac{6a}{5}$. So $DE = CH= \frac{6a}{5}$. Let the intersection of $AE$ and $BD$ be $I$. Let $AI = b$. So $IE = 18 - b$. Obviously $\triangle IDE \sim \angle IAE$ by $AA$ similarity. Since $AE = 10$ this means $IE = \sqrt{b^2 - 100}$. Now, since $\angle EAB = \angle IAB = 90^{\circ}$, we know $\triangle IAB$ is a right triangle which means $\triangle EAB \sim \triangle IAE$. Recall $\triangle AGB$ is a right triangle with legs $a$ and $15$. Thus, we have $3$ similar triangles with similarity proportions. First, we look at $\triangle ABE$ and $\triangle IAE$. We have

 $\frac{\sqrt{b^2 - 100}}{10} = \frac{b}{\sqrt{a^2 + 225}} \implies (b^2 - 100)(a^2 + 225) = 100b^2$. 

 Now, look at $\triangle IDE$ and $\triangle IAE$. We have

 $\frac{\frac{6a}{5}}{10} = \frac{18 - b}{\sqrt{b^2 - 100}} \implies 9a^2(b^2 - 100) = 625(18 - b)^2 \implies 9a^2b^2 - 900a^2 = 625(18 - b)^2$. 

 Now, recall $(b^2 - 100)(a^2 + 225) = 100b^2 \implies a^2b^2 - 100a^2 + 225b^2 - 150^2 = 100b^2 \implies 9a^2b^2 - 900a^2 = 9(100b^2 + 150^2 - 225b^2) = 9(150^2 - 125b^2)$. We now have

 $625(18 - b)^2 = 9 \cdot 25(25 \cdot 36 - 5b^2) \implies 8100 - 45b^2 = 25(18 - b)^2 \implies 1620 - 9b^2 = 1620 - 180b + 5b^2$. This is nice since we now have

 $14b^2 = 180b \implies b = 0, \frac{90}{7}$. Obviously $b = \frac{90}{7}$. Now, let's go back to 

 $(b^2 - 100)(a^2 + 225) = 100b^2 \implies (\frac{90^2}{49} - \frac{100 \cdot 49}{49})(a^2 + 225) = 100 \cdot \frac{90^2}{49}$. This means

 $(\frac{90^2 - 70^2}{49})(a^2 + 225) = 100 \cdot \frac{90^2}{49} \implies (\frac{160 \cdot 20}{49})(a^2 + 225) = 100 \cdot \frac{90^2}{49} \implies a^2 + 225 = 100 \cdot \frac{90^2}{49} \cdot \frac{49}{160 \cdot 20} = \frac{8100}{32}$. Thus, $AB = \sqrt{a^2 + 15^2} = \sqrt{a^2 + 225} = \sqrt{\frac{8100}{32}} = \frac{90}{4\sqrt{2}}$. Also, $a^2 = \frac{8100}{32} - 225 = \frac{8100 - 225 \cdot 32}{32} = \frac{25 \cdot 4 \cdot 81 - 25 \cdot 9 \cdot 32}{32} = \frac{25(324 - 288)}{32} = \frac{25 \cdot 36}{32} \implies a = \frac{30}{4\sqrt{2}}$. Finally, the area of $ABCD$ is 

 $\frac{AB + CD}{2} \cdot 18 = (\frac{90}{4\sqrt{2}} + EH + DE + HC) \cdot 9$. Now note $EH = AB = \frac{90}{4\sqrt{2}}$ and $DE = HC = \frac{6a}{5} = \frac{6}{5} \cdot \frac{30}{4\sqrt{2}} = \frac{9}{\sqrt{2}}$. Thus, we have

 $(\frac{90}{4\sqrt{2}} + \frac{90}{4\sqrt{2}} + 2 \cdot \frac{9}{\sqrt{2}})(9) = \frac{567\sqrt{2}}{2}$. Finally, multiplying this by $\sqrt{2}$ gives $\boxed{567}$. 

 




	\item (\textit{AIME I 2021}) Consider the sequence $(a_k)_{k\ge 1}$ of positive rational numbers defined by $a_1 = \frac{2020}{2021}$ and for $k\ge 1$, if $a_k = \frac{m}{n}$ for relatively prime positive integers $m$ and $n$, then


 
 \[a_{k+1} = \frac{m + 18}{n+19}.\]Determine the sum of all positive integers $j$ such that the rational number $a_j$ can be written in the form $\frac{t}{t+1}$ for some positive integer $t$.


 
 \textbf{Answer: }059



\textbf{Solution 1}

We know that $a_{1}=\frac{t}{t+1}$ when $t=2020$ so $1$ is a possible value of $j$. Note also that $a_{2}=\frac{2038}{2040}=\frac{1019}{1020}=\frac{t}{t+1}$ for $t=1019$. Then $a_{2+q}=\frac{1019+18q}{1020+19q}$ unless $1019+18q$ and $1020+19q$ are not relatively prime which happens when $q+1$ divides $18q+1019$ (by the Euclidean Algorithm), or $q+1$ divides $1001$. Thus, the least value of $q$ is $6$ and $j=2+6=8$. We know $a_{8}=\frac{1019+108}{1020+114}=\frac{1127}{1134}=\frac{161}{162}$. Now $a_{8+q}=\frac{161+18q}{162+19q}$ unless $18q+161$ and $19q+162$ are not relatively prime which happens the first time $q+1$ divides $18q+161$ or $q+1$ divides $143$ or $q=10$, and $j=8+10=18$. We have $a_{18}=\frac{161+180}{162+190}=\frac{341}{352}=\frac{31}{32}$. Now $a_{18+q}=\frac{31+18q}{32+19q}$ unless $18q+31$ and $19q+32$ are not relatively prime. This happens the first time $q+1$ divides $18q+31$ implying $q+1$ divides $13$, which is prime so $q=12$ and $j=18+12=30$. We have $a_{30}=\frac{31+216}{32+228}=\frac{247}{260}=\frac{19}{20}$. We have $a_{30+q}=\frac{18q+19}{19q+20}$, which is always reduced by EA, so the sum of all $j$ is $1+2+8+18+30=\boxed{059}$.

 \textbf{<u>Remark 1</u>}

 Whenever a fraction is in the form $\frac{t}{t+1}$ in lowest terms, the difference between the numerator and the denominator in the original fraction will always divide the numerator. We can model $a_j$ as $\frac{m+18k}{m+19k+1}$ (not necessarily simplified) if $a_{j-k}=\frac{m}{m+1}$ for integers $j$ and $k$. We want the least $k$ such that $\gcd(k+1,{m+18k})\neq1$. Let $d$ be a divisor of both $k+1$ and $m+18k$, then $d\mid18k+18$, so $d\mid{m-18}$. This follows from the Euclidean Algorithm.

 \textbf{<u>Remark 2</u>}

 The reason we look for the least $q$ in each case is because after that $q$ or $j$ value, the fraction will simplify to $m/n$ again and it won't follow the condition we defined. For example, in the $a_{2+q}=\frac{1019+18q}{1020+19q}$ case, after $j = 8$, the equation becomes useless because the fraction has simplified to something different, so we "switch over" to $a_{8+q}=\frac{161+18q}{162+19q}.$

 





\textbf{Solution 2 (Euclidean Algorithm and Generalization)}

Let $a_{j_1}, a_{j_2}, a_{j_3}, \ldots, a_{j_u}$ be all terms in the form $\frac{t}{t+1},$ where $j_1<j_2<j_3<\cdots<j_u,$ and $t$ is some positive integer.

 We wish to find $\sum_{i=1}^{u}{j_i}.$ Suppose $a_{j_i}=\frac{m}{m+1}$ for some positive integer $m.$ 

 <i>\textbf{To find $\boldsymbol{a_{j_{i+1}},}$ we look for the smallest positive integer $\boldsymbol{k'}$ for which 
 \[\boldsymbol{a_{j_{i+1}}=a_{j_i+k'}=\frac{m+18k'}{m+1+19k'}}\] is reducible:}</i>

 If $\frac{m+18k'}{m+1+19k'}$ is reducible, then there exists a common factor $d>1$ for $m+18k'$ and $m+1+19k'.$ By the \href{https://artofproblemsolving.com/wiki/index.php?title=Euclidean\_algorithm}{Euclidean Algorithm}, we have
 \begin{align*} d\mid m+18k' \text{ and } d\mid m+1+19k' &\implies d\mid m+18k' \text{ and } d\mid k'+1 \\ &\implies d\mid m-18 \text{ and } d\mid k'+1. \end{align*}
Since $m-18$ and $k'+1$ are not relatively prime, and $m$ is fixed, the smallest value of $k'$ such that $\frac{m+18k'}{m+1+19k'}$ is reducible occurs when $k'+1$ is the smallest prime factor of $m-18.$

 <i>\textbf{We will prove that for such value of $\boldsymbol{k',}$ the number $\boldsymbol{a_{j_{i+1}}}$ can be written in the form $\boldsymbol{\frac{t}{t+1}:}$}</i> 
 \[a_{j_{i+1}}=a_{j_i+k'}=\frac{m+18k'}{m+1+19k'}=\frac{(m-18)+18(k'+1)}{(m-18)+19(k'+1)}=\frac{\frac{m-18}{k'+1}+18}{\frac{m-18}{k'+1}+19}, \hspace{10mm} (*)\] where $t=\frac{m-18}{k'+1}+18$ must be a positive integer.

 We start with $m=2020$ and $a_{j_1}=a_1=\frac{2020}{2021},$ then find $a_{j_2}, a_{j_3}, \ldots, a_{j_u}$ by filling out the table below recursively:
 \[\begin{array}{c|c|c|c|c|c}  & & & & & \\ [-2ex] \boldsymbol{i} & \boldsymbol{m} & \boldsymbol{m-18} & \boldsymbol{k'+1} & \boldsymbol{k'} & \boldsymbol{a_{j_{i+1}} \left(\textbf{by } (*)\right)} \\ [0.5ex]  \hline   & & & & & \\ [-1.5ex]  1 & 2020 & 2002 & 2 & 1 & \hspace{4.25mm} a_2 = \frac{1019}{1020} \\ [1ex]      2 & 1019 & 1001 & 7 & 6 & \hspace{2.75mm} a_8 = \frac{161}{162} \\ [1ex]     3 & 161 & 143 & 11 & 10 & a_{18} = \frac{31}{32} \\ [1ex] 4 & 31 & 13 & 13 & 12 & a_{30} = \frac{19}{20} \\ [1ex] 5 & 19 & 1 & \text{N/A} & \text{N/A} & \text{N/A} \\ [1ex] \end{array}\]As $\left(j_1,j_2,j_3,j_4,j_5\right)=(1,2,8,18,30),$ the answer is $\sum_{i=1}^{5}{j_i}=\boxed{059}.$

 \textbf{<u>Remark</u>}

 Alternatively, from $(*)$ we can set 
 \[\frac{m+18k'}{m+1+19k'}=\frac{t}{t+1}.\]We cross-multiply, rearrange, and apply Simon's Favorite Factoring Trick to get 
 \[\left(k'+1\right)(t-18)=m-18.\]Since $k'+1\geq2,$ to find the smallest $k',$ we need $k'+1$ to be the smallest prime factor of $m-18.$ Now we continue with the last two paragraphs of the solution above.

 




	\item (\textit{AIME I 2021}) Let $ABCD$ be a cyclic quadrilateral with $AB=4,BC=5,CD=6,$ and $DA=7.$ Let $A_1$ and $C_1$ be the feet of the perpendiculars from $A$ and $C,$ respectively, to line $BD,$ and let $B_1$ and $D_1$ be the feet of the perpendiculars from $B$ and $D,$ respectively, to line $AC.$ The perimeter of $A_1B_1C_1D_1$ is $\frac mn,$ where $m$ and $n$ are relatively prime positive integers. Find $m+n.$


 
 \textbf{Answer: }301



\textbf{Solution 1 (Cyclic Quadrilaterals, Similar Triangles, and Trigonometry)}

This solution refers to the \textbf{Diagram} section.

 By the \textbf{Converse of the Inscribed Angle Theorem}, if distinct points $X$ and $Y$ lie on the same side of $\overleftrightarrow{PQ},$ and $\angle PXQ=\angle PYQ,$ then $P,Q,X,$ and $Y$ are cyclic. Therefore, quadrilaterals $ABA_1B_1,BCC_1B_1,CDC_1D_1,$ and $DAA_1D_1$ are all cyclic.

 Suppose $\overline{AC}$ and $\overline{BD}$ intersect at $E,$ and let $\angle AEB=\theta.$ It follows that $\angle CED=\theta$ and $\angle BEC=\angle DEA=180^\circ-\theta.$

 We obtain the following diagram:
                    \begin{figure}[H]
                        \centering
                        \includegraphics[width=0.5\textwidth]{Solution/AIME_2021_I_11_Solution_Figure_1.png}
                    \end{figure}
                    In every cyclic quadrilateral, each pair of opposite angles is supplementary. So, we have $\angle EA_1B_1=\angle EAB$ (both supplementary to $\angle B_1A_1B$) and $\angle EB_1A_1=\angle EBA$ (both supplementary to $\angle A_1B_1A$), from which $\triangle A_1B_1E \sim \triangle ABE$ by AA, with the ratio of similitude 
 \[\frac{A_1B_1}{AB}=\underbrace{\frac{A_1E}{AE}}_{\substack{\text{right} \\ \triangle A_1AE}}=\underbrace{\frac{B_1E}{BE}}_{\substack{\text{right} \\ \triangle B_1BE}}=\cos\theta. \hspace{15mm}(1)\]Similarly, we have $\angle EC_1D_1=\angle ECD$ (both supplementary to $\angle D_1C_1D$) and $\angle ED_1C_1=\angle EDC$ (both supplementary to $\angle C_1D_1C$), from which $\triangle C_1D_1E \sim \triangle CDE$ by AA, with the ratio of similitude 
 \[\frac{C_1D_1}{CD}=\underbrace{\frac{C_1E}{CE}}_{\substack{\text{right} \\ \triangle C_1CE}}=\underbrace{\frac{D_1E}{DE}}_{\substack{\text{right} \\ \triangle D_1DE}}=\cos\theta. \hspace{14.75mm}(2)\]We apply the Transitive Property to $(1)$ and $(2):$

 <ol style="margin-left: 1.5em;"> \item We get $\frac{B_1E}{BE}=\frac{C_1E}{CE}=\cos\theta,$ so $\triangle B_1C_1E \sim \triangle BCE$ by SAS, with the ratio of similitude 
 \[\frac{B_1C_1}{BC}=\frac{B_1E}{BE}=\frac{C_1E}{CE}=\cos\theta. \hspace{14.75mm}(3)\] 

 \item We get $\frac{D_1E}{DE}=\frac{A_1E}{AE}=\cos\theta,$ so $\triangle D_1A_1E \sim \triangle DAE$ by SAS, with the ratio of similitude 
 \[\frac{D_1A_1}{DA}=\frac{D_1E}{DE}=\frac{A_1E}{AE}=\cos\theta. \hspace{14mm}(4)\]

 </ol>From $(1),(2),(3),$ and $(4),$ the perimeter of $A_1B_1C_1D_1$ is
 \begin{align*} A_1B_1+B_1C_1+C_1D_1+D_1A_1&=AB\cos\theta+BC\cos\theta+CD\cos\theta+DA\cos\theta \\ &=(AB+BC+CD+DA)\cos\theta \\ &=22\cos\theta. &&\hspace{5mm}(\bigstar) \end{align*}
Two solutions follow from here:

 





\textbf{Solution 1.1 (Law of Cosines)}

Note that $\cos(180^\circ-\theta)=-\cos\theta$ holds for all $\theta.$ We apply the Law of Cosines to $\triangle ABE, \triangle BCE, \triangle CDE,$ and $\triangle DAE,$ respectively:
 \begin{alignat*}{12} &&&AE\textasciicircum{}2+BE\textasciicircum{}2-2\cdot AE\cdot BE\cdot\cos\angle AEB&&=AB\textasciicircum{}2&&\quad\implies\quad AE\textasciicircum{}2+BE\textasciicircum{}2-2\cdot AE\cdot BE\cdot\cos\theta&&=4\textasciicircum{}2, \hspace{15mm} &(1\star) \\ &&&BE\textasciicircum{}2+CE\textasciicircum{}2-2\cdot BE\cdot CE\cdot\cos\angle BEC&&=BC\textasciicircum{}2&&\quad\implies\quad BE\textasciicircum{}2+CE\textasciicircum{}2+2\cdot BE\cdot CE\cdot\cos\theta&&=5\textasciicircum{}2, \hspace{15mm} &(2\star) \\ &&&CE\textasciicircum{}2+DE\textasciicircum{}2-2\cdot CE\cdot DE\cdot\cos\angle CED&&=CD\textasciicircum{}2&&\quad\implies\quad CE\textasciicircum{}2+DE\textasciicircum{}2-2\cdot CE\cdot DE\cdot\cos\theta&&=6\textasciicircum{}2, \hspace{15mm} &(3\star) \\ &&&DE\textasciicircum{}2+AE\textasciicircum{}2-2\cdot DE\cdot AE\cdot\cos\angle DEA&&=DA\textasciicircum{}2&&\quad\implies\quad DE\textasciicircum{}2+AE\textasciicircum{}2+2\cdot DE\cdot AE\cdot\cos\theta&&=7\textasciicircum{}2. \hspace{15mm} &(4\star) \\ \end{alignat*}We subtract $(1\star)+(3\star)$ from $(2\star)+(4\star):$
 \begin{align*} 2\cdot AE\cdot BE\cdot\cos\theta+2\cdot BE\cdot CE\cdot\cos\theta+2\cdot CE\cdot DE\cdot\cos\theta+2\cdot DE\cdot AE\cdot\cos\theta&=22 \\ 2\cdot\cos\theta\cdot(\phantom{ }\underbrace{AE\cdot BE+BE\cdot CE+CE\cdot DE+DE\cdot AE}_{\text{Use the result from }\textbf{Remark}\text{.}}\phantom{ })&=22 \\ 2\cdot\cos\theta\cdot59&=22 \\ \cos\theta&=\frac{11}{59}. \end{align*}
Finally, substituting this result into $(\bigstar)$ gives $22\cos\theta=\frac{242}{59},$ from which the answer is $242+59=\boxed{301}.$

 





\textbf{Solution 1.2 (Area Formulas)}

Let the brackets denote areas. We find $[ABCD]$ in two different ways:

 <ol style="margin-left: 1.5em;"> \item Note that $\sin(180^\circ-\theta)=\sin\theta$ holds for all $\theta.$ By area addition, we get
 \begin{align*} [ABCD]&=[ABE]+[BCE]+[CDE]+[DAE] \\ &=\frac12\cdot AE\cdot BE\cdot\sin\angle AEB+\frac12\cdot BE\cdot CE\cdot\sin\angle BEC+\frac12\cdot CE\cdot DE\cdot\sin\angle CED+\frac12\cdot DE\cdot AE\cdot\sin\angle DEA \\ &=\frac12\cdot AE\cdot BE\cdot\sin\theta+\frac12\cdot BE\cdot CE\cdot\sin\theta+\frac12\cdot CE\cdot DE\cdot\sin\theta+\frac12\cdot DE\cdot AE\cdot\sin\theta \\ &=\frac12\cdot\sin\theta\cdot(\phantom{ }\underbrace{AE\cdot BE+BE\cdot CE+CE\cdot DE+DE\cdot AE}_{\text{Use the result from }\textbf{Remark}\text{.}}\phantom{ }) \\ &=\frac12\cdot\sin\theta\cdot59. \end{align*}
 

 \item By Brahmagupta's Formula, we get 
 \[[ABCD]=\sqrt{(s-AB)(s-BC)(s-CD)(s-DA)}=2\sqrt{210},\] where $s=\frac{AB+BC+CD+DA}{2}=11$ is the semiperimeter of $ABCD.$

 </ol>Equating the expressions for $[ABCD],$ we have
 \[\frac12\cdot\sin\theta\cdot59=2\sqrt{210},\] so $\sin\theta=\frac{4\sqrt{210}}{59}.$ Since $0^\circ<\theta<90^\circ,$ we have $\cos\theta>0.$ It follows that 
 \[\cos\theta=\sqrt{1-\sin^2\theta}=\frac{11}{59}.\]Finally, substituting this result into $(\bigstar)$ gives $22\cos\theta=\frac{242}{59},$ from which the answer is $242+59=\boxed{301}.$

 \textbf{Remark (Ptolemy's Theorem)} In $ABCD,$ we have
 \begin{align*} AE\cdot BE+BE\cdot CE+CE\cdot DE+DE\cdot AE &= (AE+CE)(BE+DE) \\ &=AC\cdot BD \\ &=AB\cdot CD+BC\cdot DA &&\hspace{10mm}\text{Ptolemy's Theorem} \\ &=59. \end{align*}
~MRENTHUSIASM

 





\textbf{Solution 2 (Finding cos x)}

The angle $\theta$ between diagonals satisfies 
 \[\tan{\frac{\theta}{2}}=\sqrt{\frac{(s-b)(s-d)}{(s-a)(s-c)}}\] (see \href{https://en.wikipedia.org/wiki/Cyclic\_quadrilateral\#Angle\_formulas}{https://en.wikipedia.org/wiki/Cyclic\_quadrilateral\#Angle\_formulas}).Thus, 
 \[\tan{\frac{\theta}{2}}=\sqrt{\frac{(11-4)(11-6)}{(11-5)(11-7)}}\text{ or }\tan{\frac{\theta}{2}}=\sqrt{\frac{(11-5)(11-7)}{(11-4)(11-6)}}.\]That is, $\tan^2{\frac{\theta}{2}}=\frac{1-\cos^2{\frac{\theta}{2}}}{\cos^2{\frac{\theta}{2}}}=\frac{24}{35}$ or $\frac{35}{24}$. Thus, $\cos^2{\frac{\theta}{2}}=\frac{35}{59}$ or $\frac{24}{59}$. So,
 \[\cos{\theta}=2\cos^2{\frac{\theta}{2}}-1=\pm\frac{11}{59}.\]In this context, $\cos{\theta}>0$. Thus, $\cos{\theta}=\frac{11}{59}$. The perimeter of $A_1B_1C_1D_1$ is
 \[22\cdot\cos{\theta}=22\cdot\frac{11}{59}=\frac{242}{59},\] and the answer is $m+n=242+59=\boxed{301}$.

 





\textbf{Solution 3 (Pythagorean Theorem)}

We assume that the two quadrilateral mentioned in the problem are similar (due to both of them being cyclic). Note that by Ptolemy’s, one of the diagonals has length $\sqrt{4 \cdot 6 + 5 \cdot 7} = \sqrt{59}.$ [I don't believe this is correct... are the two diagonals of $ABCD$ necessarily congruent? -peace09]* WLOG we focus on diagonal $BD.$ To find the diagonal of the inner quadrilateral, we drop the altitude from $A$ and $C$ and calculate the length of $A_1C_1.$ Let $x$ be $A_1D$ (Thus $A_1B = \sqrt{59} - x.$ By Pythagorean theorem, we have 
 \[49 - x^2 = 16 - \left(\sqrt{59} - x\right)^2 \implies 92 = 2\sqrt{59}x \implies x = \frac{46}{\sqrt{59}} = \frac{46\sqrt{59}}{59}.\] Now let $y$ be $C_1D.$ (thus making $C_1B = \sqrt{59} - y$). Similarly, we have 
 \[36 - y^2 = 25 - \left(\sqrt{59} - y\right)^2 \implies 70 = 2\sqrt{59}y \implies y = \frac{35}{\sqrt{59}} = \frac{35\sqrt{59}}{59}.\] We see that $A_1C_1$, the scaled down diagonal is just $x - y = \frac{11\sqrt{59}}{59},$ which is $\frac{\frac{11\sqrt{59}}{59}}{\sqrt{59}} = \frac{11}{59}$ times our original diagonal $BD,$ implying a scale factor of $\frac{11}{59}.$ Thus, due to perimeters scaling linearly, the perimeter of the new quadrilateral is simply $\frac{11}{59} \cdot 22 = \frac{242}{59},$ making our answer $242+59 = \boxed{301}.$

 
 \[\]$*$Indeed, the diagonals do not have length $\sqrt{59}$. However, here's why it works out: Let the length of diagonal BD be a. Then, by the same logic as given in the solution by fidgetboss\_4000, $x=\frac{33+a^{2}}{2a}$, $y=\frac{11+a^{2}}{2a}$, so $x-y=\frac{11}{a}$. Then, why does dividing by a to get a ratio work? It's because the orientation of quadrilaterals is different! In other words, instead of diagonal BD corresponding to diagonal $B_1$ $D_1$, it corresponds to diagonal $A_1$ $C_1$. Thus, to get the right ratio, we can find the ratio of the multiplied diagonals, then take the square root (this square root part is crucial. It happens because we essentially are taking something two-dimensional on the numerator and denominator, so to make it a one-dimensional similarity ratio, we must take the square root to half the dimension on each part). This yields $\sqrt{\frac{121}{a^{2}\cdot b^{2}}}$=$\frac{11}{a\cdot b}$. However, by Ptolemy, $a\cdot b$=59. Thus, the ratio is indeed $\frac{11}{59}$. ~MATH-TITAN

 





\textbf{Solution 4 (Symmetry)}

<i>\textbf{Solution}</i>

 
                    \begin{figure}[H]
                        \centering
                        \includegraphics[width=0.5\textwidth]{Solution/AIME_2021_I_11_Solution_Figure_2.png}
                    \end{figure}
                    In accordance with \textbf{Claim 1,} the ratios of pairs of one-color segments are the same and equal to $\cos \theta,$ where $\theta$ is the acute angle between the diagonals.
 \begin{align*}  s &= A'B' + B'C' + C'D' + D'A' \\    &= (AB + BC + CD + DA)\cos \theta \\    &= (a + b + c + d)\cos \theta \\    &= 22\cos \theta. \end{align*}
In accordance with \textbf{Claim 2,}
 \begin{align*} 2(ac + bd)\cos \theta = |d^2 – c^2 + b^2 – a^2|.\end{align*}

 \[2 \cdot 59 \cos \theta = |13 + 9|.\]
 \[s = 22\cos \theta = \frac{22 \cdot 11}{59} = \frac{242}{59}.\]Therefore, the answer is $242+59=\boxed{301}.$

 <i>\textbf{Claim 1}</i>

 
                    \begin{figure}[H]
                        \centering
                        \includegraphics[width=0.5\textwidth]{Solution/AIME_2021_I_11_Solution_Figure_3.png}
                    \end{figure}
                    In the triangle $ABC$, the points $B'$ and $C'$ are the bases of the heights dropped from the vertices $B$ and $C$, respectively.  $\angle A = \alpha$. Then
 \[B'C'=\begin{cases} BC \cos\alpha, & \textrm{if } \alpha < 90^\circ, \\ BC \cos (180^\circ – \alpha), & \textrm{if } \alpha >90^\circ. \end{cases}=|BC\cos{\alpha}|\]

 <i>\textbf{Proof}</i>

 Denote the orthocenter by  $A'$. Quadrilateral $B'C'BC$ is inscribed in a circle with diameter $BC$, so the marked $\angle B = \angle B'.$

 If  $\alpha < 90^\circ,$ the $\triangle AB'C' \sim \triangle ABC,$ the similarity coefficient is $AC' : AC = \cos \alpha.$So $B'C' : BC = \cos \alpha.$ 

 If  $\alpha > 90^\circ,$ the $\triangle A'B'C' \sim \triangle A'BC,$ the similarity coefficient is $A'C' : A'C = \cos (180^\circ – \alpha).$ So $B'C' : BC = \cos (180^\circ – \alpha).$

 <i>\textbf{Claim 2}</i>

 
                    \begin{figure}[H]
                        \centering
                        \includegraphics[width=0.5\textwidth]{Solution/AIME_2021_I_11_Solution_Figure_4.png}
                    \end{figure}
                    Given an inscribed quadrilateral $ABCD$ with sides $AB = a, BC = b, CD = c,$ and $DA = d.$ Prove that the $\angle \theta < 90^\circ$ between the diagonals is given by
 \begin{align*} \cos \theta = \frac {|d^2 – c^2 + b^2 – a^2|}{2(ac + bd)}.\end{align*}
<i>\textbf{Proof}</i>

 Let the point $B'$ be symmetric to $B$ with respect to the perpendicular bisector $AC.$ Then the quadrilateral $AB'CD$ is an inscribed one, $AB' = b, B'C = a.$

 
 \[2 \angle AEB = \overset{\Large\frown} {AB}  + \overset{\Large\frown} {CD}.\]
 \begin{align*} 2\angle B'AD = \overset{\Large\frown} {B'C} + \overset{\Large\frown} {CD} = \overset{\Large\frown} {AB} + \overset{\Large\frown} {CD} \implies \angle AEB = \angle B'AD\end{align*}


 We apply the Law of Cosines to $\triangle AB'D$ and $\triangle CB'D$:
 \begin{align*} B'D^2 = AD^2 + AB'^2 – 2 AD \cdot AB' \cos \theta,\end{align*}

 \[B'D^2 = CD^2 + CB'^2 + 2 CD \cdot CB' \cos \theta,\]
 \[d^2 + b^2 – 2 bd \cos \theta = c^2 + a^2 + 2ac \cos \theta,\]
 \[2(ac + bd) \cos \theta = |d^2 – c^2 + b^2 – a^2|.\]\textbf{vladimir.shelomovskii@gmail.com, vvsss}

 \textbf{Remark} This problem is kind of similar to \href{https://artofproblemsolving.com/wiki/index.php?title=2021\_AIME\_II\_Problems/Problem\_12}{2021 AIME II Problems/Problem 12}.

 




	\item (\textit{AIME I 2021}) Let $A_1A_2A_3\ldots A_{12}$ be a dodecagon ($12$-gon). Three frogs initially sit at $A_4,A_8,$ and $A_{12}$. At the end of each minute, simultaneously, each of the three frogs jumps to one of the two vertices adjacent to its current position, chosen randomly and independently with both choices being equally likely. All three frogs stop jumping as soon as two frogs arrive at the same vertex at the same time. The expected number of minutes until the frogs stop jumping is $\frac mn$, where $m$ and $n$ are relatively prime positive integers. Find $m+n$.


 
 \textbf{Answer: }019



\textbf{Solution 1}

Define the \textbf{distance} between two frogs as the number of sides between them that do not contain the third frog.

 Let $E(a,b,c)$ denote the expected number of minutes until the frogs stop jumping, such that the distances between the frogs are $a,b,$ and $c$ (in either clockwise or counterclockwise order). Without the loss of generality, assume that $a\leq b\leq c.$

 We wish to find $E(4,4,4).$ Note that:

 <ol style="margin-left: 1.5em;"> \item At any moment before the frogs stop jumping, the only possibilities for $(a,b,c)$ are $(4,4,4),(2,4,6),$ and $(2,2,8).$ 

 \item $E(a,b,c)$ does not depend on the actual positions of the frogs, but depends on the distances between the frogs. 

 \item At the end of each minute, each frog has $2$ outcomes. So, there are $2^3=8$ outcomes in total.

 </ol>We have the following system of equations:
 \begin{align*} E(4,4,4)&=1+\frac{2}{8}E(4,4,4)+\frac{6}{8}E(2,4,6), \\ E(2,4,6)&=1+\frac{4}{8}E(2,4,6)+\frac{1}{8}E(4,4,4)+\frac{1}{8}E(2,2,8), \\ E(2,2,8)&=1+\frac{2}{8}E(2,2,8)+\frac{2}{8}E(2,4,6). \end{align*}
Rearranging and simplifying each equation, we get
 \begin{align*} E(4,4,4)&=\frac{4}{3}+E(2,4,6), &(1) \\ E(2,4,6)&=2+\frac{1}{4}E(4,4,4)+\frac{1}{4}E(2,2,8), &\hspace{12.75mm}(2) \\ E(2,2,8)&=\frac{4}{3}+\frac{1}{3}E(2,4,6). &(3) \end{align*}
Substituting $(1)$ and $(3)$ into $(2),$ we obtain 
 \[E(2,4,6)=2+\frac{1}{4}\left[\frac{4}{3}+E(2,4,6)\right]+\frac{1}{4}\left[\frac{4}{3}+\frac{1}{3}E(2,4,6)\right],\] from which $E(2,4,6)=4.$ Substituting this into $(1)$ gives $E(4,4,4)=\frac{16}{3}.$

 Therefore, the answer is $16+3=\boxed{019}.$

 





\textbf{Solution 2 (Markov Chain)}


            \begin{enumerate}[label=(\roman*)]
                \item <i>Main article: \href{https://artofproblemsolving.com/wiki/index.php?title=Markov\_Chains}{Markov Chains}</i>
            \end{enumerate}
            We can solve the problem by removing $1$ frog, and calculate the expected time for the remaining $2$ frogs. In the original problem, when the movement stops, $2$ of the $3$ frogs meet. Because the $3$ frogs cannot meet at one vertex, the probability that those two specific frogs meet is $\frac13$. If the expected time for the two frog problem is $E'$, then the expected time for the original problem is $\frac 13 E'$.

 The distance between the two frogs can only be $0$, $2$, $4$, $6$. We use the distances as the states to draw the following \href{https://en.wikipedia.org/wiki/Markov\_chain}{Markov Chain}. This Markov Chain is much simpler than that of Solution 1 Supplement in the \textbf{Remark} section.

 
                    \begin{figure}[H]
                        \centering
                        \includegraphics[width=0.5\textwidth]{Solution/AIME_2021_I_12_Solution_Figure_1.png}
                    \end{figure}
                    
 \begin{align*} E(2) &= 1 + \frac12 \cdot E(2) + \frac14 \cdot E(4)\\ E(4) &= 1 + \frac14 \cdot E(2) + \frac12 \cdot E(4) + \frac14 \cdot E(6)\\ E(6) &= 1 + \frac12 \cdot E(4) + \frac12 \cdot E(6) \end{align*}


 By solving the above system of equations, $E(4) = 16$. The answer for the original problem is $\frac{16}{3}$, $16 + 3 = \boxed{\textbf{019}}$

 


 





\textbf{Solution 3 (Pure algebraic manipulation, rigorous)}

Let $(x,y,z)$ be the number of positions between every pair of adjacent frogs. Notice that each of these three numbers are always odd as the number of positions can only change by $+2$, $+0$ or $-2$ in every move. We also have that $x+y+z=9$ as the frog themselves occupy $3$ positions.

 
Let $a_n$ to be the probability that the frogs are at $(3,3,3)$ Situation at the $n$th move.

 Let $b_n$ to be the probability that the frogs are at $(1,3,5)$ Situation at the $n$th move.

 Let $c_n$ to be the probability that the frogs are at $(1,1,7)$ Situation at the $n$th move.

 These are all the cases with $x,y,z$ all odd and sum to $9$.

 
After listing all the $8$ movements, we get the recurrence

 $a_{n+1}=\frac{1}{4}a_n+\frac{1}{8}b_n$

 $b_{n+1}=\frac{3}{4}a_n+\frac{1}{2}b_n+\frac{1}{4}c_n$

 $c_{n+1}=\frac{1}{8}b_n+\frac{1}{4}c_n$

 
A $3,3,3$ Situation never leads to the end.

 A $1,3,5$ Situation has $\frac{1}{4}$ chance of leading to the end.

 A $1,1,7$ Situation has $\frac{1}{2}$ chance of leading to the end.

 
Then we calculate the probability of the process ending at the $n$th move.

 
$P($ending at move $1)$$=0$

 $P($ending at move $2)$$=\frac{1}{4}b_1+\frac{1}{2}c_1=\frac{3}{16}$

 $P($ending at move $3)$$=\frac{1}{4}b_2+\frac{1}{2}c_2=\frac{3}{16}$

 $P($ending at move $4)$$=\frac{1}{4}b_3+\frac{1}{2}c_3$

 $=\frac{1}{4}(\frac{3}{4}a_2+\frac{1}{2}b_2+\frac{1}{4}c_2)+\frac{1}{2}(\frac{1}{8}b_2+\frac{1}{4}c_2)$ (by plugging in the equations above)

 $=\frac{3}{16}a_2+\frac{1}{8}b_2+\frac{1}{16}c_2+\frac{1}{16}b_2+\frac{1}{8}c_2$

 $=\frac{3}{16}a_2+\frac{3}{16}b_2+\frac{3}{16}c_2$

 $=\frac{3}{16}(a_2+b_2+c_2)$

 $a_2+b_2+c_2$ is the probability of the frogs being at a non-ending situation at the $2$nd move. This is equivalent to the process not ending on or before the $2$nd move.

 The probability here can be calculated reversely, it is equal to $1-$ (the probability of the process ending on or before the $2$nd move).

 It is equal to $1-$ (the probability of the process ending at the $1$st move) $-$ (the probability of the process ending at the $2$nd move).

 $P($ending at move $4)$$=\frac{3}{16}(1-\frac{3}{16})=\frac{39}{256}$

 
Then continue calculations like this.

 
$P($ending at move $5)$$=\frac{1}{4}b_4+\frac{1}{2}c_4$

 $=\frac{3}{16}(a_3+b_3+c_3)$

 $=\frac{3}{16}$$P($the process not ending on or before the $3$rd move$)$

 $=\frac{3}{16}$$(1-P($the process ending on or before the $3$rd move$))$

 $=\frac{3}{16}$$(1-P($the process ending on the $1$st move$)$$-P($the process ending on the $2$nd move$)$$-P($the process ending on the $3$rd move$))$

 $=\frac{3}{16}(1-\frac{3}{16}-\frac{3}{16})$

 $=\frac{15}{128}$

 
$P($ending at move $6)$$=\frac{1}{4}b_5+\frac{1}{2}c_5$

 $=\frac{3}{16}(a_4+b_4+c_4)$

 $=\frac{3}{16}$$P($the process not ending on or before the $4$th move$)$

 $=\frac{3}{16}$$(1-P($the process ending on or before the $4$th move$))$

 $=\frac{3}{16}(1-\frac{3}{16}-\frac{3}{16}-\frac{39}{256})$

 $=\frac{363}{4096}$

 
.........

 
In general, $P($ending at move ${n+1})$$=\frac{1}{4}b_n+\frac{1}{2}c_n$

 $=\frac{3}{16}(a_{n-1}+b_{n-1}+c_{n-1})$

 $=\frac{3}{16}$$P($the process not ending on or before the $n-1$th move$)$

 $=\frac{3}{16}$$(1-P($the process ending on or before the $n-1$th move$))$

 $=\frac{3}{16}$$(1-\sum_{i=1}^{n-1}P($the process ending at the $i$th move$))$

 Notice that $P($ending at move ${n+2})$ has just one more term to minus than $P($ending at move ${n+1})$. Subtracting the equation for $n+2$ and the equation for $n+1$, we derive

 $P($ending at move ${n+1})-$$P($ending at move ${n+2})$$=\frac{3}{16}P($ending at move $n)$

 
Write $P($ending at move $n)$ as $Q_n$.

 We have $Q_{n+1}-Q_{n+2}=\frac{3}{16}Q_n$.

 We have $Q_{n+2}=Q_{n+1}-\frac{3}{16}Q_n$

 
We can actually prove the convergence of the expected value now in order to be rigorous. $Q_n$ is a probability, a non-negative real number. For all postive integers $n$ we have $Q_{n+2}\leq Q_{n+1}$, the sequence is non-increasing since $Q_3$.

 $Q_{n+2}=Q_{n+1}-\frac{3}{16}Q_n\leq Q_n-\frac{3}{16}Q_n=\frac{13}{16}Q_n$ for $n\geq3$.

 Thus each term of the sequence $Q$ is no more than the corresponding term of two alternating geometric sequences, one for its odd terms and one for its even terms. The infinite sum of $Q$ converges, and the expected value $Q_1+2Q_2+3Q_3+4Q_4+......$ converges as well.

 
Now we calculate $Q_1+Q_2+Q_3+Q_4+Q_5......$. We set this finite number equal to $S$.

 $S=Q_2+Q_3+Q_4+Q_5+......$ (because $Q_1$ is $0$)

 $S=\frac{3}{16}+Q_2-\frac{3}{16}Q_1+Q_3-\frac{3}{16}Q_2+Q_4-\frac{3}{16}Q_3+Q_5-\frac{3}{16}Q_4+......$

 $S=\frac{3}{16}+S-\frac{3}{16}S$

 $\frac{3}{16}-\frac{3}{16}S=0$

 $S=1$

 
Knowing $S=1$, we can find the expected value $Q_1+2Q_2+3Q_3+4Q_4+5Q_5+......$.

 
Set this finite value equal to $E$.

 $E=Q_1+2Q_2+3Q_3+4Q_4+5Q_5+......$

 $E=\frac{3}{8}+3(Q_2-\frac{3}{16}Q_1)+4(Q_3-\frac{3}{16}Q_2)+5(Q_4-\frac{3}{16}Q_3)+6(Q_5-\frac{3}{16}Q_4)+......$

 $E=\frac{3}{8}+3Q_2+4Q_3+5Q_4+6Q_5+......-\frac{3}{16}(3Q_1+4Q_2+5Q_3+6Q_4+......)$

 $E=\frac{3}{8}+E+S-\frac{3}{16}(4Q_2+5Q_3+6Q_4+7Q_5......)$

 $E=\frac{3}{8}+E+S-\frac{3}{16}(E+2S)$

 $E=\frac{3}{8}+E+1-\frac{3}{16}(E+2)$

 $3(E+2)=22$

 $E=\frac{16}{3}$

 
The answer to the problem is $16+3=19$.

 \textbf{Remark (Markov Chain)} 





\textbf{Solution 1 Supplement}

Solution 1 can be represented by the following \href{https://en.wikipedia.org/wiki/Markov\_chain}{Markov Chain}:

 
                    \begin{figure}[H]
                        \centering
                        \includegraphics[width=0.5\textwidth]{Solution/AIME_2021_I_12_Solution_Figure_2.png}
                    \end{figure}
                    
                \begin{itemize}
                    \item From state $(4, 4, 4)$ to state $(4, 4, 4)$: the $3$ frogs must jump in the same direction, $2 \cdot \frac18 = \frac14$.
                \end{itemize}


                
                \begin{itemize}
                    \item From state $(4, 4, 4)$ to state $(2, 4, 6)$: $2$ frogs must jump in the same direction, and the other must jump in the opposite direction, $\binom32 \cdot 2 \cdot \frac18 = \frac34$.
                \end{itemize}


                
                \begin{itemize}
                    \item From state $(2, 4, 6)$ to state $(4, 4, 4)$: the $2$ frogs with a distance of $4$ in between must jump in the same direction so that they will be further away from the other frog, and the other frog must jump in the opposite direction as those $2$ frogs, $\frac18$.
                \end{itemize}


                
                \begin{itemize}
                    \item From state $(2, 4, 6)$ to state $(2, 4, 6)$: the $3$ frogs can all jump in the same direction; or the $2$ frogs with a distance of $2$ in between jumps away from each other and the other frog jumps closer to the closest frog; or the $2$ frogs with a distance of $6$ in between jump closer to each other and away from the third frog, the third frog jumps closer to the closest frog; $(2 + 1 + 1) \cdot \frac18 = \frac12$.
                \end{itemize}


                
                \begin{itemize}
                    \item From state $(2, 4, 6)$ to state $(2, 2, 8)$: the $2$ frogs with a distance of $2$ in between must jump closer to the other frog and the other frog must jump close to those $2$ frogs, $\frac18$.
                \end{itemize}


                
                \begin{itemize}
                    \item From state $(2, 2, 8)$ to state $(2, 4, 6)$: $2$ frogs that have no frogs in between must both jump in the same direction away from the other frog, the other frog must also jump away from those $2$ frogs, $2 \cdot \frac18 = \frac14$.
                \end{itemize}


                
                \begin{itemize}
                    \item From state $(2, 2, 8)$ to state $(2, 2, 8)$: the $3$ frogs must all jump in the same direction, $2 \cdot \frac18 = \frac14$.
                \end{itemize}


                
                \begin{itemize}
                    \item From state $(2, 2, 8)$ to state $(0, x, y)$: frogs with a distance of $2$ must jump closer to each other, the other frog can jump in any direction, $\frac12 \cdot \frac12 \cdot 2 = \frac12$.
                \end{itemize}


                
                \begin{itemize}
                    \item From state $(2, 4, 6)$ to state $(0, x, y)$: the frogs with a distance of $2$ must jump closer to each other, the other frog can jump in any direction, $\frac12 \cdot \frac12 = \frac14$.
                \end{itemize}


                Because $a + b + c = 12$, we can use $(a, b)$ to represent the states which is simpler than using $(a, b, c)$ in Solution 1.

 \textbf{Summary} The two above Markov Chains are \href{https://en.wikipedia.org/wiki/Absorbing\_Markov\_chain}{Absorbing Markov Chain}. The state of $2$ frogs meeting is the absorbing state. This problem asks for the \href{https://en.wikipedia.org/wiki/Absorbing\_Markov\_chain\#Expected\_number\_of\_steps}{Expected Number of Steps} before being absorbed in the absorbing state.

 Let $p_{ij} = P(X_{n+1} = j | X_n = i)$, the probability that state $i$ transits to state $j$ on the next step. 

 Let $e_i$ be the expected number of steps before being absorbed in the absorbing state when starting from transient state $i$.

 $e_i = \sum_{j} [p_{ij} \cdot ( 1 + e_{j})] = \sum_{j} (p_{ij} + p_{ij} \cdot e_{j}) = \sum_{j} p_{ij} + \sum_{j} (p_{ij} \cdot e_{j}) = 1 + \sum_{j} (p_{ij} \cdot e_{j})$

 $e_i$ is $1$ plus the sum of the products of $p_{ij}$ and $s_j$ of all the next state $j$, $\sum_{j} p_{ij} = 1$.

 \href{https://artofproblemsolving.com/wiki/index.php?title=2014\_AMC\_12B\_Problems/Problem\_22}{2014 AMC 12B Problem 22} and \href{https://artofproblemsolving.com/wiki/index.php?title=2017\_AMC\_12A\_Problems/Problem\_22}{2017 AMC 12A Problem 22} are similar problem with simpler states, both problems can be solved by Absorbing Markov Chain.

 




	\item (\textit{AIME I 2021}) Circles $\omega_1$ and $\omega_2$ with radii $961$ and $625$, respectively, intersect at distinct points $A$ and $B$. A third circle $\omega$ is externally tangent to both $\omega_1$ and $\omega_2$. Suppose line $AB$ intersects $\omega$ at two points $P$ and $Q$ such that the measure of minor arc $\widehat{PQ}$ is $120^{\circ}$. Find the distance between the centers of $\omega_1$ and $\omega_2$.


 
 \textbf{Answer: }672



\textbf{Solution 1 (Radical Axis)}

Let $O_i$ and $r_i$ be the center and radius of $\omega_i$, and let $O$ and $r$ be the center and radius of $\omega$.

 Since $\overline{AB}$ extends to an arc with arc $120^\circ$, the distance from $O$ to $\overline{AB}$ is $r/2$. Let $X=\overline{AB}\cap \overline{O_1O_2}$. Consider $\triangle OO_1O_2$. The line $\overline{AB}$ is perpendicular to $\overline{O_1O_2}$ and passes through $X$. Let $H$ be the foot from $O$ to $\overline{O_1O_2}$; so $HX=r/2$. We have by tangency $OO_1=r+r_1$ and $OO_2=r+r_2$. Let $O_1O_2=d$.
                    \begin{figure}[H]
                        \centering
                        \includegraphics[width=0.5\textwidth]{Solution/AIME_2021_I_13_Solution_Figure_1.png}
                    \end{figure}
                    Since $X$ is on the radical axis of $\omega_1$ and $\omega_2$, it has equal power with respect to both circles, so
 \[O_1X^2 - r_1^2 = O_2X^2-r_2^2 \implies O_1X-O_2X = \frac{r_1^2-r_2^2}{d}\]since $O_1X+O_2X=d$. Now we can solve for $O_1X$ and $O_2X$, and in particular,
 \begin{align*} O_1H &= O_1X - HX = \frac{d+\frac{r_1^2-r_2^2}{d}}{2} - \frac{r}{2} \\ O_2H &= O_2X + HX = \frac{d-\frac{r_1^2-r_2^2}{d}}{2} + \frac{r}{2}.  \end{align*}
We want to solve for $d$. By the Pythagorean Theorem (twice):
 \begin{align*} &\qquad -OH^2 = O_2H^2 - (r+r_2)^2 = O_1H^2 - (r+r_1)^2 \\ &\implies \left(d+r-\frac{r_1^2-r_2^2}{d}\right)^2 - 4(r+r_2)^2 = \left(d-r+\frac{r_1^2-r_2^2}{d}\right)^2 - 4(r+r_1)^2 \\ &\implies 2dr - 2(r_1^2-r_2^2)-8rr_2-4r_2^2 = -2dr+2(r_1^2-r_2^2)-8rr_1-4r_1^2 \\ &\implies 4dr = 8rr_2-8rr_1 \\ &\implies d=2r_2-2r_1 \end{align*}
Therefore, $d=2(r_2-r_1) = 2(961-625)=\boxed{672}$.

 





\textbf{Solution 2 (Linearity)}

Let $O_{1}$, $O_{2}$, and $O$ be the centers of $\omega_{1}$, $\omega_{2}$, and $\omega$ with $r_{1}$, $r_{2}$, and $r$ their radii, respectively. Then, the distance from $O$ to the radical axis $\ell\equiv\overline{AB}$ of $\omega_{1}, \omega_{2}$ is equal to $\frac{1}{2}r$. Let $x=O_{1}O_{2}$ and $O^{\prime}$ the orthogonal projection of $O$ onto line $\ell$. Define the function $f:\mathbb{R}^{2}\rightarrow\mathbb{R}$ by 
 \[f(X)=\text{Pow}_{\omega_{1}}(X)-\text{Pow}_{\omega_{2}}(X).\] Then 
 \begin{align*} f(O_{1})=-r_{1}^{2}-(x-r_{2})(x+r_{2})&=-x^{2}+r_{2}^{2}-r_{1}^{2}, \\ f(O_{2})=(x-r_{1})(x+r_{1})-(-r_{2}^{2})&=x^{2}+r_{2}^{2}-r_{1}^{2}, \\ f(O)=r(r+2r_{1})-r(r+2r_{2})&=2r(r_{1}-r_{2}), \\ f(O^{\prime})&=0. \end{align*}
 By \href{https://services.artofproblemsolving.com/download.php?id=YXR0YWNobWVudHMvMC84LzkzZjZjZmFlMGViY2E3MDMxNWQzY2IzNzFlZTk5NWFmOTM5ZGY1LnBkZg==&rn=TGluZWFyaXR5IG9mIFBvd2VyIG9mIGEgUG9pbnQucGRm}{linearity}, 
 \[\frac{f(O_{2})-f(O_{1})}{f(O)-f(O^{\prime})}=\frac{O_{1}O_{2}}{OO^{\prime}}=\frac{x}{\frac{1}{2}r}=\frac{2x}{r}.\] Notice that $f(O_{2})-f(O_{1})=x^{2}-(-x^{2})=2x^{2}$ and $f(O)-f(O^{\prime})=2r(r_{1}-r_{2})$, thus 
 \begin{align*}\frac{2x^{2}}{2r(r_{1}-r_{2})}&=\frac{2x}{r}\end{align*}
 Dividing both sides by $\frac{2x}{r}$ (which is obviously nonzero as $x$ is nonzero) gives us 
 \begin{align*}\frac{x}{2(r_{1}-r_{2})}&=1\end{align*}
 so $x=2(r_{1}-r_{2})$. Since $r_{1}=961$ and $r_{2}=625$, the answer is $x=2\cdot(961-625)=\boxed{672}$.

 





\textbf{Solution 3}

Denote by $O_1$, $O_2$, and $O$ the centers of $\omega_1$, $\omega_2$, and $\omega$, respectively. Let $R_1 = 961$ and $R_2 = 625$ denote the radii of $\omega_1$ and $\omega_2$ respectively, $r$ be the radius of $\omega$, and $\ell$ the distance from $O$ to the line $AB$. We claim that
 \[\dfrac{\ell}{r} = \dfrac{R_2-R_1}{d},\]where $d = O_1O_2$. This solves the problem, for then the $\widehat{PQ} = 120^\circ$ condition implies $\frac{\ell}r = \cos 60^\circ = \frac{1}{2}$, and then we can solve to get $d = \boxed{672}$.
                    \begin{figure}[H]
                        \centering
                        \includegraphics[width=0.5\textwidth]{Solution/AIME_2021_I_13_Solution_Figure_2.png}
                    \end{figure}
                    

 Denote by $O_1$ and $O_2$ the centers of $\omega_1$ and $\omega_2$ respectively. Set $X$ as the projection of $O$ onto $O_1O_2$, and denote by $Y$ the intersection of $AB$ with $O_1O_2$. Note that $\ell = XY$. Now recall that
 \[d(O_2Y-O_1Y) = O_2Y^2 - O_1Y^2 = R_2^2 - R_1^2.\]Furthermore, note that
 \begin{align*}d(O_2X - O_1X) &= O_2X^2 - O_1X^2= O_2O^2 - O_1O^2 \\ &= (R_2 + r)^2 - (R_1+r)^2 = (R_2^2 - R_1^2) + 2r(R_2 - R_1).\end{align*}
Substituting the first equality into the second one and subtracting yields
 \[2r(R_2 - R_1) = d(O_2X - O_1X) - d(O_2Y - O_1Y) = 2dXY,\]which rearranges to the desired.

 





\textbf{Solution 4 (Quick)}

Suppose we label the points as shown \href{https://services.artofproblemsolving.com/download.php?id=YXR0YWNobWVudHMvNC9mLzRiM2JjYThjYmZlY2ViZGI0ODhjYzE4YzMyMmM0M2QyOTZlMmU5LmpwZw==&rn=MTU4ODUxMDg3XzczMDI0ODE4MTAwNjA5N184NDQzMjQxMjM3MDQ2NzQ5NjM4X24uanBn}{below}. 
                    \begin{figure}[H]
                        \centering
                        \includegraphics[width=0.5\textwidth]{Solution/AIME_2021_I_13_Solution_Figure_3.png}
                    \end{figure}
                    By radical axis, the tangents to $\omega$ at $D$ and $E$ intersect on $AB$. Thus $PDQE$ is harmonic, so the tangents to $\omega$ at $P$ and $Q$ intersect at $X \in DE$. Moreover, $OX \parallel O_1O_2$ because both $OX$ and $O_1O_2$ are perpendicular to $AB$, and $OX = 2OP$ because $\angle POQ = 120^{\circ}$. Thus
 \[O_1O_2 = O_1Y - O_2Y = 2 \cdot 961 - 2\cdot 625 = \boxed{672}\]by similar triangles.

 





\textbf{Solution 5 (Official MAA)}

Like in other solutions, let $O$ be the center of $\omega$ with $r$ its radius; also, let $O_{1}$ and $O_{2}$ be the centers of $\omega_{1}$ and $\omega_{2}$ with $R_{1}$ and $R_{2}$ their radii, respectively. Let line $OP$ intersect line $O_{1}O_{2}$ at $T$, and let $u=TO_{2}$, $v=TO_{1}$, $x=PT$, where the length $O_{1}O_{2}$ splits as $u+v$. Because the lines $PQ$ and $O_{1}O_{2}$ are perpendicular, lines $OT$ and $O_{1}O_{2}$ meet at a $60^{\circ}$ angle.

 Applying the Law of Cosines to $\triangle O_{2}PT$, $\triangle O_{1}PT$, $\triangle O_{2}OT$, and $\triangle O_{1}OT$ gives 
 \begin{align*}\triangle O_{2}PT&:O_{2}P^{2}=u^{2}+x^{2}-ux \\ \triangle O_{1}PT&:O_{1}P^{2}=v^{2}+x^{2}+vx \\ \triangle O_{2}OT&:(r+R_{2})^{2}=u^{2}+(r+x)^{2}-u(r+x) \\ \triangle O_{1}OT&:(r+R_{1})^{2}=v^{2}+(r+x)^{2}+v(r+x)\end{align*}


 Adding the first and fourth equations, then subtracting the second and third equations gives us 
 \[\left(O_{2}P^{2}-O_{1}P^{2}\right)+\left(R_{1}^{2}-R_{2}^{2}\right)+2r(R_{1}-R_{2})=r(u+v)\]

 Since $P$ lies on the radical axis of $\omega_{1}$ and $\omega_{2}$, the power of point $P$ with respect to either circle is 
 \[O_{2}P^{2}-R_{2}^{2}=O_{1}P^{2}-R_{1}^{2}.\]

 Hence $2r(R_{1}-R_{2})=r(u+v)$ which simplifies to 
 \[u+v=2(R_{1}-R_{2}).\]

 The requested distance 
 \[O_{1}O_{2}=O_{1}T+O_{2}T=u+v\] is therefore equal to $2\cdot(961-625)=\boxed{672}$.

 





\textbf{Solution 6 (Geometry)}


                    \begin{figure}[H]
                        \centering
                        \includegraphics[width=0.5\textwidth]{Solution/AIME_2021_I_13_Solution_Figure_4.png}
                    \end{figure}
                    Let circle $\omega$ tangent circles $\omega_1$ and $\omega_2,$ respectively at distinct points $C$ and $D$. Let $O, O_1, O_2 (r, r_1, r_2)$ be the centers (the radii) of $\omega, \omega_1$ and $\omega_2,$ respectively. WLOG $r_1 < r_2.$ Let $F$ be the point of $OO_2$ such, that $OO_1 =OF.$ Let $M$ be the midpoint $FO_1, OE \perp AB, CT$ be the radical axes of  $\omega_1$ and $\omega, T \in AB.$ 

 Then $T$ is radical center of  $\omega, \omega_1$ and $\omega_2, TD = CT.$

 In $\triangle OFO_1 (OF = OO_1) OT$ is bisector of $\angle O, OM$ is median 

 $\hspace{10mm} \implies O,T,$  and $M$ are collinear. 

 
 \[\angle OCT = \angle ODT = \angle OET = 90^\circ \implies\]

 $OCTDE$ is cyclic (in $\Omega), OT$ is diameter $\Omega.$ $O_1O_2 \perp AB, OM \perp FO_1 \implies \angle FO_1O_2 = \angle OTE$$\angle OTE = \angle ODE$ as they intercept the same arc  $\overset{\Large\frown}{OE}$ in $\Omega.$
 \[OE \perp AB, O_1O_2 \perp AB \implies O_1 O_2 || OE \implies\]
 \[\angle OO_2O_1 = \angle O_2 OE \implies \triangle ODE \sim \triangle O_2 O_1 F \implies\]
 \[\frac {OE}{OD} = \frac {O_2F}{O_1O_2} \implies \cos \frac {120^\circ}{2} = \frac{r_2 + r - r_1 -r} {O_1O_2}\implies {O_1O_2}= 2|r_2 – r_1|.\]

 Since $r_{1}=625$ and $r_{2}=961$, the answer is $2\cdot|961-625|=\boxed{672}$.

 \textbf{vladimir.shelomovskii@gmail.com, vvsss}

 





\textbf{Solution 7}

We are not given the radius of circle $w$, but based on the problem statement, that radius isn't important. We can set $w$ to have radius infinity (solution 8), but if you didn't observe that, you could also set the radius to be $2r$ so that the line containing the center of $w$, call it $W$, and $w_2$, call it $W_2$, is perpendicular to the line containing the center of $w_1$, call it $W_1$ and $w_2$. Let $AB = 2h$ and $W_1W_2 = x$. Also, let the projections of $W$ and $W_1$ onto line $AB$ be $X$ and $Y$, respectively.

 By \href{https://artofproblemsolving.com/wiki/index.php?title=Pythagorean\_Theorem}{Pythagorean Theorem} on $\triangle WW_1W_2$, we get
 \[x^2+(625+2r)^2=(961+2r)^2 \;(1)\]Note that since $\angle PWQ = 120$, $\angle PWX = 60$. So, $WX = 2r/2 = r = W_1Y$. We now get two more equations from Pythag:
 \[h^2+r^2 = 625^2 \; (2)\]
 \[h^2+(x-r)^2 = 961^2 \; (3)\]From subtracting $(2)$ and $(3)$, $x^2-2rx=961^2-625^2 \; (4)$. Rearranging $(1)$ yields $x^2-1344r = 961^2-625^2$. Plugging in our result from $(4)$, $x^2-1344r= x^2-2rx \implies 1344r = 2rx \implies x=\boxed{672}$.

 





\textbf{Solution 8 (Cheese)}

Let the \href{https://artofproblemsolving.com/wiki/index.php?title=Circle}{circle} $\omega$ be infinitely big (a line). Then for it to be split into an \href{https://artofproblemsolving.com/wiki/index.php?title=Arc}{arc} of $120^{\circ}$, $\overline{PQ}$ must intersect at a $60^{\circ}$ with line $\omega$.

 Notice the 30-60-90 triangle in the image. $O_1R = 961 - 625$.

 Thus, the distance between the centers of $\omega_1$ and $\omega_2$ is $2(961-625)=\boxed{672}$

 
                    \begin{figure}[H]
                        \centering
                        \includegraphics[width=0.5\textwidth]{Solution/AIME_2021_I_13_Solution_Figure_5.png}
                    \end{figure}
                    


 


 


 


 


 


 


 


 


 


 


 


 


 


 


 


 


 


 


 


 


 
picture by afly

 Who wanted to see animated video solutions can see this. I found this really helpful. 

 P.S: This video is not made by me. And solution is same like below solutions.

 ≈@rounak138

 




	\item (\textit{AIME I 2021}) For any positive integer $a,$ $\sigma(a)$ denotes the sum of the positive integer divisors of $a$. Let $n$ be the least positive integer such that $\sigma(a^n)-1$ is divisible by $2021$ for all positive integers $a$. Find the sum of the prime factors in the prime factorization of $n$.


 
 \textbf{Answer: }125



\textbf{Solution 1}

We first claim that $n$ must be divisible by $42$. Since $\sigma(a^n)-1$ is divisible by $2021$ for all positive integers $a$, we can first consider the special case where $a$ is prime and $a \neq 0,1 \pmod{43}$. By Dirichlet's Theorem (Refer to the \textbf{Remark} section.), such $a$ always exists.

 Then $\sigma(a^n)-1 = \sum_{i=1}^n a^i = a\left(\frac{a^n - 1}{a-1}\right)$. In order for this expression to be divisible by $2021=43\cdot 47$, a necessary condition is $a^n - 1 \equiv 0 \pmod{43}$. By \href{https://artofproblemsolving.com/wiki/index.php?title=Fermat\%27s\_Little\_Theorem}{Fermat's Little Theorem}, $a^{42} \equiv 1 \pmod{43}$. Moreover, if $a$ is a primitive root modulo $43$, then $\text{ord}_{43}(a) = 42$, so $n$ must be divisible by $42$.

 By similar reasoning, $n$ must be divisible by $46$, by considering $a \not\equiv 0,1 \pmod{47}$.

 We next claim that $n$ must be divisible by $43$. By Dirichlet, let $a$ be a prime that is congruent to $1 \pmod{43}$. Then $\sigma(a^n) \equiv n+1 \pmod{43}$, so since $\sigma(a^n)-1$ is divisible by $43$, $n$ is a multiple of $43$.

 Alternatively, since $\left(\frac{a(a^n - 1^n)}{a-1}\right)$ must be divisible by $43,$ by LTE, we have $v_{43}(a)+v_{43}{(a-1)}+v_{43}{(n)}-v_{43}{(a-1)} \geq 1,$ which simplifies to $v_{43}(n) \geq 1,$ which implies the desired result.

 Similarly, $n$ is a multiple of $47$.

 Lastly, we claim that if $n = \text{lcm}(42, 46, 43, 47)$, then $\sigma(a^n) - 1$ is divisible by $2021$ for all positive integers $a$. The claim is trivially true for $a=1$ so suppose $a>1$. Let $a = p_1^{e_1}\ldots p_k^{e_k}$ be the prime factorization of $a$. Since $\sigma(n)$ is \href{https://artofproblemsolving.com/wiki/index.php?title=Multiplicative\_function}{multiplicative}, we have
 \[\sigma(a^n) - 1 = \prod_{i=1}^k \sigma(p_i^{e_in}) - 1.\]We can show that $\sigma(p_i^{e_in}) \equiv 1 \pmod{2021}$ for all primes $p_i$ and integers $e_i \ge 1$, so
 \[\sigma(p_i^{e_in}) = 1 + (p_i + p_i^2 + \ldots + p_i^n) + (p_i^{n+1} + \ldots + p_i^{2n}) + \ldots + (p_i^{n(e_i-1)+1} + \ldots + p_i^{e_in}),\]where each expression in parentheses contains $n$ terms. It is easy to verify that if $p_i = 43$ or $p_i = 47$ then $\sigma(p_i^{e_in}) \equiv 1 \pmod{2021}$ for this choice of $n$, so suppose $p_i \not\equiv 0 \pmod{43}$ and $p_i \not\equiv 0 \pmod{47}$. Each expression in parentheses equals $\frac{p_i^n - 1}{p_i - 1}$ multiplied by some power of $p_i$. If $p_i \not\equiv 1 \pmod{43}$, then FLT implies $p_i^n - 1 \equiv 0 \pmod{43}$, and if $p_i \equiv 1 \pmod{43}$, then $p_i + p_i^2 + \ldots + p_i^n \equiv 1 + 1 + \ldots + 1 \equiv 0 \pmod{43}$ (since $n$ is also a multiple of $43$, by definition). Similarly, we can show $\sigma(p_i^{e_in}) \equiv 1 \pmod{47}$, and a simple \href{https://artofproblemsolving.com/wiki/index.php?title=Chinese\_Remainder\_Theorem}{CRT} argument shows $\sigma(p_i^{e_in}) \equiv 1 \pmod{2021}$. Then $\sigma(a^n) \equiv 1^k \equiv 1 \pmod{2021}$.

 Then the prime factors of $n$ are $2,3,7,23,43,$ and $47,$ and the answer is $2+3+7+23+43+47 = \boxed{125}$.

 





\textbf{Solution 2 (Along the Lines of Solution 1)}

$n$ only needs to satisfy $\sigma(a^n)\equiv 1 \pmod{43}$ and $\sigma(a^n)\equiv 1 \pmod{47}$ for all $a$. Let's work on the first requirement (mod 43) first. All $n$ works for $a=1$. If $a>1$, let $a$'s prime factorization be $a=p_1^{e_1}p_2^{e_2}\cdots p_k^{e_k}$. The following three statements are the same:

 <ol style="margin-left: 1.5em;"> \item For all $a$, we have $\sigma(a^n)=\prod_{i=1}^k(1+p_i+\cdots+p_i^{ne_i})\equiv1\pmod{43}$. 

 \item For all $e>0$ and prime $p$, we have $1+p+\cdots+p^{ne}\equiv1\pmod{43}$. 

 \item For all prime $p$, we have $p+p^2+\cdots+p^n \equiv 0 \pmod{43}$.

 </ol>We can show this by casework on $p$:

 <ol style="margin-left: 1.5em;"> \item For $p\equiv0\pmod{43}$, no matter what $n$ is, it is true that 
 \[p+p^2+\cdots+p^n \equiv 0 \pmod{43}.\] 

 \item For all $p\equiv1\pmod{43}$, it is true that 
 \[p+p^2+\cdots+p^n \equiv n \equiv 0 \pmod{43}.\] One can either use brute force or Dirichlet's Theorem to show such $p$ exists. Therefore, $43|n$. 

 \item For all $p\not\equiv0,1\pmod{43}$, it is true that
 \[p+p^2+\cdots+p^n \equiv 0 \pmod{43} \Leftrightarrow p^n-1\equiv0\pmod{43}.\] According to Fermat's Little Theorem, $42|n$ is sufficient. To show it's necessary, we just need to show $43$ has a prime primitive root. This can be done either by brute force or as follows. $43$ is prime and it must have a primitive root $t\neq 1$ that's coprime to $43$. All numbers of the form $43k+t$ are also primitive roots of $43$. According to Dirichlet's Theorem there must be infinitely many primes among them.

 </ol>Similar arguments for modulo $47$ lead to $46|n$ and $47|n$. Therefore, we get $n=\operatorname{lcm}[42,43,46,47]$. Following the last paragraph of Solution 1 gives the answer $\boxed{125}$.

 





\textbf{Solution 3 (Casework)}

We perform casework on $a:$

 <ol style="margin-left: 1.5em;"> <li>$a=1.$ For all positive integers $n,$ we have $\sigma\left(a^n\right)-1=0.$

 </li> <li>$a$ is a prime number. For all prime numbers $p,$ note that
 \[\sigma\left(p^n\right)-1=\left(\sum_{i=0}^{n}p^i\right)-1=\sum_{i=1}^{n}p^i=\frac{p^{n+1}-p}{p-1}=\frac{p\left(p^n-1\right)}{p-1}\]by geometric series. 

 Recall that $2021=43\cdot47.$ Applying the Chinese Remainder Theorem, we get the following system of congruences:
 \begin{align*} \frac{p\left(p^n-1\right)}{p-1}&\equiv0\pmod{43}, &(1)\\ \frac{p\left(p^n-1\right)}{p-1}&\equiv0\pmod{47}. &(2) \end{align*}
We perform casework on $(1):$

 <ol style="margin-left: 1.5em;" type="A"> <li>$p-1\not\equiv0\pmod{43}.$ We need $p\left(p^n-1\right)\equiv0\pmod{43}.$ 

 If $p=43,$ then this congruence is true for all positive integers $n.$ 

 If $p\neq43,$ then this congruence becomes $p^n-1\equiv0\pmod{43}.$ By Fermat's Little Theorem, we obtain $n\equiv0\pmod{42}.$

 </li> <li>$p-1\equiv0\pmod{43}.$ We need $p^n-1$ to contain more factors of $43$ than $p-1$ does. More generally, for all positive integers $k,$ if $p-1\equiv0 \ \left(\operatorname{mod} \ 43^k\right),$ then $p^n-1\equiv0 \ \left(\operatorname{mod} \ 43^{k+1}\right).$ 

 Let $p=43^km+1$ for some positive integer $m$ such that $m\not\equiv0\pmod{43}.$ We substitute this into $p^n-1\equiv0 \ \left(\operatorname{mod} \ 43^{k+1}\right)$ and apply the Binomial Theorem:
 \begin{align*} \left(43^km+1\right)^n-1&\equiv0 \ \left(\operatorname{mod} \ 43^{k+1}\right) \\ \Biggl[\binom{n}{0}\left(43^km\right)^0+\binom{n}{1}\left(43^km\right)^1+\phantom{ }\underbrace{\binom{n}{2}\left(43^km\right)^2+\cdots+\binom{n}{n}\left(43^km\right)^n}_{0 \ \left(\operatorname{mod} \ 43^{k+1}\right)}\phantom{ }\Biggr]-1 &\equiv0 \ \left(\operatorname{mod} \ 43^{k+1}\right) \\ \left[1+43^kmn\right]-1&\equiv0 \ \left(\operatorname{mod} \ 43^{k+1}\right) \\ 43^kmn&\equiv0 \ \left(\operatorname{mod} \ 43^{k+1}\right), \end{align*}
from which $n\equiv0\pmod{43}.$

 </li></ol>For $(1),$ we find that $n\equiv0\pmod{42}$ and $n\equiv0\pmod{43}.$ For $(2),$ we find that $n\equiv0\pmod{46}$ and $n\equiv0\pmod{47}$ by a similar argument. 

 Together, we conclude that $n=\operatorname{lcm}(42,43,46,47)$ is the least such positive integer $n$ for this case.

 </li> 

 <li>$a$ is a composite number. Let $a=\prod_{i=1}^{u}p_i^{e_i}$ be the prime factorization of $a.$ Note that $\sigma(a)$ is a multiplicative function: 
 \[\sigma(a)=\sigma\left(\prod_{i=1}^{u}p_i^{e_i}\right)=\prod_{i=1}^{u}\sigma\left(p_i^{e_i}\right)=\prod_{i=1}^{u}\left(\sum_{j=0}^{e_i}p_i^j\right),\] as this product of geometric series spans all positive divisors of $a.$ 

 At $n=\operatorname{lcm}(42,43,46,47),$ it follows that
 \begin{align*} \sigma\left(a^n\right)-1&=\left(\prod_{i=1}^{u}\sigma\left(p_i^{e_in}\right)\right)-1 \\ &\equiv\left(\prod_{i=1}^{u}1\right)-1 &&\pmod{2021} \\ &\equiv0 &&\pmod{2021}. \end{align*}


 </li>

 </ol>Finally, the least such positive integer $n$ for all cases is
 \begin{align*} n&=\operatorname{lcm}(42,43,46,47) \\ &=\operatorname{lcm}(2\cdot3\cdot7,43,2\cdot23,47) \\ &=2\cdot3\cdot7\cdot23\cdot43\cdot47, \end{align*}
so the sum of its prime factors is $2+3+7+23+43+47=\boxed{125}.$

 





\textbf{Solution 4 (Cheap)}

Since the problem works for all positive integers $a$, let's plug in $a=2$ and see what we get. Since $\sigma({2^n}) = 2^{n+1}-1,$ we have $2^{n+1} \equiv 2 \pmod{2021}.$ Simplifying using CRT and \href{https://artofproblemsolving.com/wiki/index.php?title=Fermat\%27s\_Little\_Theorem}{Fermat's Little Theorem}, we get that $n \equiv 0 \pmod{42}$ and $n \equiv 0 \pmod{46}.$ Then, we can look at $a$ being a $1\pmod{43}$ prime and a $1\pmod{47}$ prime, just like in Solution 1, to find that $43$ and $47$ also divide $n$. There don't seem to be any other odd "numbers" to check, so we can hopefully assume that the answer is the sum of the prime factors of $\text{lcm(42, 43, 46, 47)}.$ From here, follow solution 1 to get the final answer $\boxed{125}$.

 





\textbf{Solution 5 (Similar to Solution 4 and USEMO 2019 Problem 4)}

Warning: This solution doesn't explain why $43*47\mid n$.

 It says the problem implies that it works for all positive integers $a$, we basically know that If $p \equiv 1 \pmod q$, then from "USEMO 2019 Problem 4", if $p^n + p^{n-1} + \cdots + 1 \equiv 1 \pmod{q}$, then 
 \[\frac{p^{en+1}-1}{p-1} = p^{en} + p^{en-1} + \cdots + 1 \equiv 1 \pmod{q}.\] From here we can just let $\sigma(2^n)$ or be a power of $2$ which we can do 
 \[\sigma(2^n)=1+2+2^2+2^3+2^4+\cdots+2^n=2^{n+1}-1,\] which is a geometric series. We can plug in $a=2$ like in Solution 4 and use CRT. We have the prime factorization $2021 = 43 \cdot 47$. We use CRT to find that 
 \begin{align*} 2^n &\equiv 1 \pmod{43}, \\ 2^n &\equiv 1 \pmod{47}. \end{align*}
 We see that this is just FLT which is $a^{p-1} \equiv 1 \pmod p$ we see that all multiples of $42$ will work for first and $46$ for the second. We can figure out that it is just $\text{lcm}(43-1,47-1)\cdot43\cdot47$ which when we add up we get that it's just the sum of the prime factors of $\text{lcm}(42,43,46,47)$ which you can just look at Solution 1 to find out the sum of the prime factors and get the answer $\boxed{125}$.

 





\textbf{Solution 6 (Easier than most solutions)}

The sum of the factors of a number with prime factorization $n = p_1^{a_1} p_2^{a_2} \cdots p_k^{a_k}$ is given by:
 \[\sigma(n) = \prod_{i=1}^k \left( \sum_{j=0}^{a_i} p_i^j \right) = \prod_{i=1}^k \frac{p_i^{a_i+1} - 1}{p_i - 1}\],or in more simplicity, The sum of the factors of a number with prime factorization \( (a_1^{e_1} a_2^{e_2} \cdots a_k^{e_k}) \) is given by:
 \[\sigma(n) = (1 + a_1 + a_1^2 + \cdots + a_1^{e_1}) (1 + a_2 + a_2^2 + \cdots + a_2^{e_2}) \cdots (1 + a_k + a_k^2 + \cdots + a_k^{e_k})\]

 Then, notice that if $a$ is prime then $\sigma(a^n) - 1 \equiv 0 \pmod{2021}$, then $(\sum_{i=0}^{n} a^i) - 1 \equiv 0 \pmod{2021}$, or $\sum_{i=1}^{n} a^i \equiv 0 \pmod{2021}$, which is to say $a\left(\frac{a^n - 1}{a - 1}\right) \equiv 0 \pmod{2021}$, and is thus equivalent to $0 \pmod{43}~\text{and}~ 0 \pmod{47}$.

 If $a \equiv 0 \pmod{43}$ then note that $\frac{a^n - 1}{a - 1} \equiv 1 \pmod{43}$, in which a$^n \equiv a \pmod{43}$ and thus $a^{n-1} \equiv 1 \pmod{43}$. Then, because 43 is prime, we have that $n$ must be a multiple of 42. In addition, $n$ must be a multiple of 46 through the same logic.

 Then, think about when $a$ is composite. Then, it may be written as $(\sum_{i=1}^{n})^k p^{i}_1 \cdot p^{i}_2 \cdot ... \cdot p^{i}_k$ (note that $(\sum_{i=1}^{n})^k$ implies $k$ nested sums) where each $p_k$ is a prime. This, when subtracted by 1, must be equivalent to $0 \pmod{2021}$ and thus each individual sum is the sum of a geometric series with every value being equivalent to some $0 \pmod{2021}$. 

 Now, if we take the expanded form we essentially find
 \[\sigma(n) = (1 + a_1 + a_1^2 + \cdots + a_1^{e_1}) (1 + a_2 + a_2^2 + \cdots + a_2^{e_2}) \cdots (1 + a_k + a_k^2 + \cdots + a_k^{e_k}) - 1 \equiv 0 \pmod {2021}\]or 
 \[\sigma(n) = (1 + a_1 + a_1^2 + \cdots + a_1^{e_1}) (1 + a_2 + a_2^2 + \cdots + a_2^{e_2}) \cdots (1 + a_k + a_k^2 + \cdots + a_k^{e_k}) \equiv 1 \pmod {2021}\]

 Which then implies at least one of these sums is equivalent to $1\pmod{43}$ or/and $1 \pmod{47}$. Then, WLOG, say that sum is $1 + a_1 + a_1^2 + \cdots + a_1^{e_1}$ for $\mod 43$.

 Then, $\frac{a_1^{e_1+1}-1}{a_1 - 1}$, in which $a_1^{e_1} \equiv 1 \pmod{43}$. The only way this is possible is if $e_1 = 42$. This then implies the primilaritveness of $e_1$, as thus $\sigma(a^n) \equiv n+1 \pmod{43}$. Then, $n$ is divisible by 43, and the same logic applies for 47.

 Another way to approach that is to realize that once $e_1 = 42$, the factor of $n$ must be divisible by 43 through the sum of factors formula.

 The smallest $n$ is then the $lcm(42,43,46,47)$, and continue as in solution 1.

 \textbf{Remark (Dirichlet's Theorem)} All solutions require Dirichlet's Theorem, which states that for any coprime positive integers $k$ and $r$, there are infinitely many primes $p$ congruent to $r\pmod{k}$.

 




	\item (\textit{AIME I 2021}) Let $S$ be the set of positive integers $k$ such that the two parabolas
 \[y=x^2-k~~\text{and}~~x=2(y-20)^2-k\]intersect in four distinct points, and these four points lie on a circle with radius at most $21$. Find the sum of the least element of $S$ and the greatest element of $S$.


 
 \textbf{Answer: }285



\textbf{Solution 1 (Inequalities and Circles)}

Note that $y=x^2-k$ is an upward-opening parabola with the vertex at $(0,-k),$ and $x=2(y-20)^2-k$ is a rightward-opening parabola with the vertex at $(-k,20).$ We consider each condition separately:

 <ol style="margin-left: 1.5em;"> \item The two parabolas intersect at four distinct points.

 By a quick sketch, we have two subconditions:<ol style="margin-left: 1.5em;" type="A"> <li>The point $(-k,20)$ is on or below the parabola $y=x^2-k.$We need $20\leq(-k)^2-k,$ from which $k\geq5.$ 

 Moreover, the point $(-k,20)$ is on the parabola $y=x^2-k$ when $k=5.$ We will prove that the two parabolas intersect at four distinct points at this value of $k:$

 Substituting $y=x^2-5$ into $x=2(y-20)^2-5,$ we get $x=2\left(\left(x^2-5\right)-20\right)^2-5.$ Expanding and rearranging give 
 \[2x^4-100x^2-x+1245=0. \hspace{20mm}(\bigstar)\]By either the graphs of the parabolas or the Rational Root Theorem, we conclude that $x=-5$ is a root of $(\bigstar).$ So, we factor its left side: 
 \[(x+5)\left(2x^3-10x^2-50x+249\right)=0.\]By either the graphs of the parabolas or Descartes' Rule of Signs, we conclude that $2x^3-10x^2-50x+249=0$ has two positive roots and one negative root such that $x\neq-5.$ So, $(\bigstar)$ has four distinct real roots, or the two parabolas intersect at four distinct points.

 For Subcondition A, we deduce that $k\geq5.$ 

 <u>Remark for Subcondition A</u>

 Recall that if $1\leq k\leq 4,$ then the point $(-k,20)$ is above the parabola $y=x^2-k.$ It follows that for $-k\leq x\leq0:$

 <ul style="list-style-type:square;">\item The maximum value of $y$ for the parabola $y=x^2-k$ occurs at $x=-k,$ from which $y=k^2-k\leq12.$

 \item The minimum value of $y$ for the parabola $x=2(y-20)^2-k$ occurs at $x=0,$ from which $y=20-\sqrt{\frac k2}>18.$

 </ul>Clearly, the parabola $x=2(y-20)^2-k$ and the left half of the parabola $y=x^2-k$ do not intersect. Therefore, the two parabolas do not intersect at four distinct points.

 </li> \item The point $(0,-k)$ is on or below the parabola $x=2(y-20)^2-k.$The lower half of the parabola $x=2(y-20)^2-k$ is $y=20-\sqrt{\frac{x+k}{2}}.$ We need $-k\leq20-\sqrt{\frac k2},$ which holds for all values of $k.$

 For Subcondition B, we deduce that $k$ can be any positive integer.

 </ol>\textbf{For Condition 1, we obtain $\boldsymbol{k\geq5}$ by taking the intersection of Subconditions A and B.}

  \item The four points of intersection lie on a circle with radius at most $21.$For equations of circles, the coefficients of $x^2$ and $y^2$ must be the same. So, we add the equation $y=x^2-k$ to half the equation $x=2(y-20)^2-k:$ 
 \[y+\frac12x=x^2+(y-20)^2-\frac32k.\]We expand, rearrange, and complete the squares:
 \begin{align*} y+\frac12x&=x^2+y^2-40y+400-\frac32k \\ \frac32k-400&=\left(x^2-\frac12x\right)+\left(y^2-41y\right) \\ \frac32k-400+\frac{1}{16}+\frac{1681}{4}&=\left(x-\frac14\right)^2+\left(y-\frac{41}{2}\right)^2. \end{align*}
We need $\frac32k-400+\frac{1}{16}+\frac{1681}{4}\leq21^2,$ from which $k\leq\left\lfloor\frac{6731}{24}\right\rfloor=280.$

 \textbf{For Condition 2, we obtain $\boldsymbol{k\leq280.}$}

 </ol>Taking the intersection of Conditions 1 and 2 produces $5\leq k\leq280.$ Therefore, the answer is $5+280=\boxed{285}.$

 





\textbf{Solution 2 (Translations, Inequalities, Circles)}

Make the translation $y \rightarrow y+20$ to obtain $20+y=x^2-k$ and $x=2y^2-k$. Multiply the first equation by $2$ and sum, we see that $2(x^2+y^2)=3k+40+2y+x$. Completing the square gives us $\left(y- \frac{1}{2}\right)^2+\left(x - \frac{1}{4}\right)^2 = \frac{325+24k}{16}$; this explains why the two parabolas intersect at four points that lie on a circle*. For the upper bound, observe that $LHS \leq 21^2=441 \rightarrow 24k \leq 6731$, so $k \leq 280$. 

 For the lower bound, we need to ensure there are $4$ intersections to begin with. (Here I'm using the un-translated coordinates.) Draw up a graph, and realize that two intersections are guaranteed, on the so called "right branch" of $y=x^2-k$. As we increase the value of $k$, two more intersections appear on the "left branch":

 $k=4$ does not work because the "leftmost" point of $x=2(y-20)^2-4$ is $(-4,20)$ which lies to the right of $\left(-\sqrt{24}, 20\right)$, which is on the graph $y=x^2-4$. While technically speaking this doesn't prove that there are no intersections (why?), drawing the graph should convince you that this is the case. Clearly, $k<4$ does not work.

 $k=5$ does work because the two graphs intersect at $(-5,20)$, and by drawing the graph, you realize this is not a tangent point and there is in fact another intersection nearby, due to slope. Therefore, the answer is $5+280=\boxed{285}$.

 
                \begin{itemize}
                    \item In general (assuming four intersections exist), when two conics intersect, if one conic can be written as $ax^2+by^2=f(x,y)$ and the other as $cx^2+dy^2=g(x,y)$ for polynomials $f$ and $g$ of degree at most $1$, whenever $(a,b),(c,d)$ are linearly independent (L.I.), we can combine the two equations and then complete the square to achieve $(x-p)^2+(y-q)^2=r^2$. We can also combine these two equations to form a parabola, or a hyperbola, or an ellipse. When $(a,b),(c,d)$ are not L.I., the intersection points instead lie on a line, which is a circle of radius infinity. When the two conics only have $3,2$ or $1$ intersection point(s), the statement that all these points lie on a circle is trivially true.
                \end{itemize}


                





\textbf{Solution 3 (Parabola's Properties)}

<i>\textbf{Claim}</i>

 
                    \begin{figure}[H]
                        \centering
                        \includegraphics[width=0.5\textwidth]{Solution/AIME_2021_I_15_Solution_Figure_1.png}
                    \end{figure}
                    Let the axes of two parabolas be perpendicular, their focal parameters be  $p_1$ and $p_2$ and the distances from the foci to the point of intersection of the axes be  $x_2$ and $y_1$. Suppose that these parabolas intersect at four points.

 Then these points lie on the circle centered at point $(p_2, p_1)$ with radius $r = \sqrt{2(p_1^2 + p_2^2 + p_1 y_1 + p_2 x_2)}.$

 <i>\textbf{Proof}</i>

 Let's introduce a coordinate system with the center at the point of intersection of the axes. Let the first (red) parabola have axis $x = 0,$ focal parameter  $p_1$ and focus at point $A(0, –y_1), y_1 > 0.$ Let second (blue) parabola have axis $y = 0,$ focal parameter $p_2$ and focus at point $B(–x_2,0), x_2  > 0.$Let us denote the angle between the vector connecting the focus of the first parabola and its point and the positive direction of the ordinate axis $2\theta,$ its length $\rho_1(\theta),$ the angle between the vector connecting the focus of the second parabola and its point and the positive direction of the abscissa axis $2\phi,$ its length $\rho_2(\phi).$ Then
 \[\rho_1(\theta) = \frac{p_1}{1 - \cos(2\theta)}, \rho_2(\phi) = \frac{p_2}{1 - \cos(2\phi)}.\]

 Abscissa of the point of intersection is
 \begin{align*} x =\rho_1 \sin(2\theta) = p_1\cot\theta = \rho_2 \cos (2\phi) - x_2 = \frac{p_2}{2} (\cot^2\phi - 1)- x_2,\end{align*}

 \begin{align*} x^2 = p_1^2 \cot ^2 \theta , 2 p_1\cot\theta = p_2 \cos^2 \phi - p_2 - 2x_2 .\end{align*}
Ordinate of the point of intersection is
 \begin{align*} y =\rho_2 \sin 2\phi = p_2\cot\phi = \rho_1 \cos 2\theta - y_1 = \frac{p_1}{2} (\cot^2\theta - 1)- y_1,\end{align*}

 \begin{align*} y^2 = p_2^2 \cot ^2 \phi , 2 p_2\cot\phi = p_1 \cos^2 \theta - p_1 - 2y_1 .\end{align*}
The square of the distance from point of intersection to the point $(p_2, p_1)$ is 
 \begin{align*} r^2 = (x-p_2)^2 + (y-p_1)^2 = x^2 + y^2 - 2 p_1 y - 2 p_2 x + p_1^2 + p_2^2 .\end{align*}
After simple transformations, we get $r^2 = 2(p_1^2 + p_2^2 + p_1 y_1 + p_2 x_2).$

 Hence, any intersection point has the same distance $r$ from the point  $(p_2, p_1).$

 
<i>\textbf{Solution}</i>

 Parameters of the parabola $y = x^2 – k$ are $p_1 = \frac{1}{2}, y_1 = 20 + k – \frac{1}{2}.$

 Parameters of the parabola $\frac{x}{2} = (y – 20)^2 – \frac{k}{2}$ are $p_2 = \frac{1}{4}, x_2 = k – \frac{1}{4}  \implies r^2 = 20 + \frac{3k}{2}.$

 If $r \le 21, k \le \frac{842}{3},$ then integer  $k \le 280.$

 The vertex of the second parabola is point  $(– k,20)$  can be on the parabola $y = x^2 – k$ or below the point of the parabola with the same abscissa. So  
 \[20 \ge (– k)^2 – k  \implies   5 \le k \le 280.\] Therefore, the answer is $5+280=\boxed{285}$.

 \textbf{vladimir.shelomovskii@gmail.com, vvsss}

 





\textbf{Solution 4}


                    \begin{figure}[H]
                        \centering
                        \includegraphics[width=0.5\textwidth]{Solution/AIME_2021_I_15_Solution_Figure_2.png}
                    \end{figure}
                    $\textbf{Min}:$

 When $k = 4$:

 $20 = x^2-4$ $x=-\sqrt{24}$

 $-\sqrt{24}<-4$ $\times$ 

 When $k = 5$:

 $20 = x^2-5$ $x=-5$

 $-5 = -5$ $\checkmark$ 

 $\textbf{Max}:$

 \begin{align*}\{ & \left.\begin{aligned}y &= x^2-k \\x &= 2y^2-80y+800-k\end{aligned} \right.\end{align*}


 
$x = 2y^2+2y-82y+800-k$

 $x = 2y^2+2x^2-2k-82y+800-k$

 $2x^2+2y^2-x-82y+800-3k=0$

 $x^2+y^2-\frac{x}{2}-41y+400-\frac{3}{2}k = 0$

 $(x-\frac{1}{4})^2+(y-\frac{41}{2})^2+400-(\frac{1}{4})^2-(\frac{41}{2})^2 - \frac{3}{2}k = 0$

 $(x-\frac{1}{4})^2+(y-\frac{41}{2})^2 = (\frac{1}{4})^2+(\frac{41}{2})^2 + \frac{3}{2}k - 400 \leq 441$

 $\frac{1}{16}+\frac{1681}{4}+\frac{3}{2}k \leq 841$

 $1+6724+24k \leq 13456$

 $24k \leq 6631$

 $k \leq \frac{6631}{24}, \quad k_{max} = 280$

 Therefore the answer is $280+5 = \boxed{285}$

 -cassphe

 The provided solution was incorrect because of a calculation error; I have corrected the answer and fixed latex typos.

 




\end{enumerate}